% 第一章 预备知识: 集合论与范畴 --- 中译
% 原书 pp. 1--40 (Aluffi, Algebra: Chapter 0, Chapter I).
% 本文件进度: §\ref{sec:prelim-naive}--§\ref{sec:universal-properties} 已译 (2026-08-28).

\chapter{预备知识:集合论与范畴}
\label{chap:prelim}


% ===== 记号索引 v1 集中埋点 (页码参考原书书页; 首次出现位置待人工校对) =====
% 本章记号条目（11 条）。OCR 已做最佳校正，但部分符号字形与首次
% 出现位置仍待逐条人工复核（页码沿用原书第 699 页 Index of Notation）。
% ===== /记号索引 =====
\index[sym1]{$An$@$An$}
\index[sym1]{$\Aut_c(A)$@$\Aut_c(A)$}
\index[sym1]{$B^A$@$B^A$}
\index[sym1]{$c$@$c$}
\index[sym1]{$\End_c(A)$@$\End_c(A)$}
\index[sym1]{$\Hom_c(A,B)$@$\Hom_c(A,B)$}
\index[sym1]{$\N$@$\N$}
\index[sym1]{$\mathcal{P}(S)$@$\mathcal{P}(S)$}
\index[sym1]{$\Q$@$\Q$}
\index[sym1]{$\R$@$\R$}
\index[sym1]{$\Z$@$\Z$}
\index{Zermelo@Zermelo!Ernst@Ernst}集合论本身就是数学的一个领域, 而对它的严格处理 (比如说, 通过著名的 Zermelo--Fränkel 公理)
远远超出本书的范围, 也超出笔者的能力. 我们只打算使用所谓的``朴素''集合论%
\index{pusujihelun@朴素集合论(naive set theory)}
, 它不过是一套记号与术语系统, 使我们能够精确地表述数学的定义、命题及其证明.

\index{daishu@代数(algebra)!jiaohuan@交换(commutative)}熟悉这套语言对于学习代数%
\index{daishu@代数(algebra)}
\index{fanchou@范畴(category)!Ab@Ab}这样的学科是必不可少的, 的确, 我们假定读者此前已经接触过它. 在本章中, 我们先回顾朴素集合论语言的一部分, 主要目的是确立全书将使用的记号; 随后我们会对范畴%
\index{fanchou@范畴(category)}
\index{daishu@代数(algebra)!chu@除(division)}的语言略作品尝, 它在代数以及许多其他领域中发挥着强大的统一作用. 我们的主要目标是传达``万有性质''%
\index{wanyouxingzhi@万有性质(universal property)}
的观念, 它将是贯穿全书的一条主线.

\section{朴素集合论}
\label{sec:prelim-naive}

\subsection{集合}
\label{subsec:sets}

集合%
\index{jih@集合(set)}
\index{yuan@元(element)}的概念将``对象的汇集''这一直观想法形式化. 一个集合由它所包含的元素%
\index{yuansu@元素(element)}
确定: 两个集合 $A$, $B$ 相等 (记作 $A = B$) 当且仅当它们包含的元素恰好相同. ``元素是什么?''在朴素集合论中是一个被禁止的问题: 追问必须止于某处. 我们可以方便地假装有一个由元素构成的``宇宙''可供使用, 并从这个宇宙中取材, 构造出我们需要的元素与集合, 同时默认我们将要探索的一切运算都能在这个宇宙中进行. (这正是微妙之处!) 无论如何, 我们通过给出一个确切的规则来确定一个集合, 该规则指明哪些元素在其中. 这样的定义通常写在花括号之间, 也可以是一个元素明确的完整清单:
\[
	A := \{1,2,3\}
\]
\footnote{``:='' 这一记号常用来表示: 左边的符号由右边的内容来定义. 逻辑上这只是在表达两侧的相等, 完全可以直接写``=''; 这个多出来的``:''不过是从计算机科学继承下来的一种心理上方便的装饰.}%
这是由整数 $1$、$2$、$3$ 构成的集合. 按照约定, 列出元素时的次序%
\footnote{有序的清单用圆括号表示: $(1,2,3)$ 与 $(1,3,2)$ 不相同.}%
以及清单中的重复都与定义无关. 于是, 同一个集合可以有许多种写法:
\[
	\{1,2,3\} = \{1,3,2\} = \{1,2,1,3,3,2,3,1,1,2,1,3\}.
\]
\index{yuan@元(element)!Abel fanchouduixiangde@Abel 范畴对象的(of an object of an abelian category)!chudeng@初等(elementary)}这种表示集合的方式可能相当笨拙, 而且无论如何只对有限集真正有效. 对于无限集, 一种流行的变通办法是写出一个清单, 其中某些元素被理解为按某个模式延伸---例如, 偶整数%
\index{ouzhengshu@偶整数(even integer)}
构成的集合可以写成
\[
	E = \{\dots, -2, 0, 2, 4, 6, \dots\},
\]
\index{ji@积(product)!zhi@直(direct)}但这样的定义天然带有歧义, 因而留下误解的空间. 此外, 有些集合简直``太大''了, 即便原则上也无法列出: 例如 (但愿读者在高等微积分中学到过) 实数实在太多, 无法像``列出''整数那样把它们一一``列出''.

更可取的做法往往是: 把集合的元素表达为某个更大的 (并且已知的) 集合 $S$ 中满足某个性质 $P$ 的那些元素 $s$. 这时可以写
\[
	A = \{s \in S \mid s \text{ 满足 } P\}
\]
($\in$ 的意思是``是\ldots 的元素\ldots''), 一般说来这是精确而无歧义的%
\footnote{但要注意, 存在罗素悖论%
\index{luosubeilun@罗素悖论(Russell's paradox)}
这样的病态情形, 它们表明即便这种定义方式也可能导致无意义的结果. 只要 $S$ 本身确实已知是一个集合, 就一切平安无事.}%
.

\index{duozhongji@多重集(multiset)!Sylow dingli@Sylow 定理(Sylow's theorem)}我们偶尔会遇到集合概念的一个变体, 称为``多重集''%
\index{duochongji@多重集(multiset)}
\index{duozhongji@多重集(multiset)!dingli@定理(theorem)}\index{duozhongji@多重集(multiset)!zhangliang@张量(tensor)}\index{duozhongji@多重集(multiset)!Tarski@Tarski!Alfred@Alfred}\index{duozhongji@多重集(multiset)!duichenxing@对称性(symmetry)}\index{duozhongji@多重集(multiset)!duichen@对称(symmetric)}\index{chongshu@重数(multiplicity)!tezhengzhide@特征值的(of an eigenvalue)!jihe@几何(geometry)}\index{chongshu@重数(multiplicity)!tezhengzhide@特征值的(of an eigenvalue)!daishu@代数(algebra)}\index{chongshu@重数(multiplicity)!duoxiangshigende@多项式根的(of a root of a polynomial)}\index{chongshu@重数(multiplicity)}. 多重集是允许元素``带重数地''出现的集合: 也就是说, 在多重集中 $\{2,2\}$ 与 $\{2\}$ 是不同的. 定义多重集的正确方式要借助映射, 我们很快就会遇到 (见例~\ref{subsec:multisets-indexed}).

几个著名的集合:
\begin{itemize}
	\item $\emptyset$: 空集%
\index{kongji@空集(empty set)}
\index{yuan@元(element)!yinzi@因子(divisor)}, 不含任何元素;
	\item $\N$: 自然数集 (即非负整数集);\index[sym1]{42@$\N$}
	\item $\Z$: 整数集;\index[sym1]{68@$\Z$}
	\item $\Q$: 有理数集;\index[sym1]{48@$\Q$}
	\item $\R$: 实数集;\index[sym1]{52@$\R$}
	\item $\CC$: 复数集.\index[sym1]{07@$\CC$}
\end{itemize}
\index{yuan@元(element)!yunsuan@运算(operation)}此外, ``单元素集''%
\index{danyuansuji@单元素集(singleton)}
一词用来指恰由一个元素构成的集合. 于是 $\{1\}$, $\{2\}$, $\{3\}$ 是互不相同的集合, 但它们都是单元素集.

下面是几个有用的符号 (称为量词%
\index{liangci@量词(quantifier)}
):
\begin{itemize}
	\item $\exists$ 表示``存在\ldots'' (存在量词);
	\item $\forall$ 表示``对所有\ldots'' (全称量词).
\end{itemize}
另外, $\exists!$ 用来表示``存在唯一的\ldots''.

例如, 偶整数集可以写成
\[
	E = \{a \in \Z \mid (\exists n \in \Z)\ a = 2n\}:
\]
\index{zi@字(word)}用文字说, ``所有使得存在整数 $n$ 满足 $a = 2n$ 的整数 $a$''. 在这个例子里, 我们可以把 $\exists$ 换成 $\exists!$ 而不改变这个集合---但这与 $\Z$ 的性质有关, 与数学句法无关. 人们还常采用简写
\[
	E = \{2n \mid n \in \Z\},
\]
其中存在量词被默认省略.

能够毫不费力地解析这样的符号串, 并且能够流利地把它们写出来, 是极其重要的. 本书假定读者已经掌握了这一技能.

注意, 书写事物的次序可能带来巨大的差别. 例如, 命题
\[
	(\forall a \in \Z) (\exists b \in \Z) \quad b = 2a
\]
为真: 它说, 把任意整数加倍的结果是整数; 而
\[
	(\exists b \in \Z) (\forall a \in \Z) \quad b = 2a
\]
为假: 它说, 存在一个固定的整数 $b$, 它``同时''是每个整数的两倍---根本没有这种东西.

还要注意, 单独写下
\[
	b = 2a
\]
并不能传达足够的信息, 除非上下文把附着于 $a$ 与 $b$ 的量词交代得完全清楚: 诚如方才所见, 不同的量词可以把同一个式子变成真命题或假命题.

\subsection{集合的包含}
\label{subsec:inclusion}

如前所述, 两个集合相等当且仅当它们包含相同的元素. 我们说集合 $S$ 是集合 $T$ 的子集%
\index{ziji@子集(subset)}
, 如果 $S$ 的每个元素都是 $T$ 的元素, 用符号写成
\[
	S \subseteq T.
\]
按照约定, $S \subset T$ 的含义相同: 也就是说 (与 $<$ 对 $\leq$ 的情形不同), 它不排除 $S$ 与 $T$ 相等的可能. 为避免任何混淆, 本书将一致地使用 $\subseteq$. 我们用 $S \subsetneq T$ 表示 $S$``真''含于 $T$: 即 $S \subseteq T$ 且 $S \neq T$.

可以用逻辑的语言来考虑``集合的包含'': $S \subseteq T$ 意味着
\[
\index{Lie@Lie!Sophus@Sophus}	s \in S \implies s \in T
\]
(量词 $\forall s$ 默认省略); 也就是说, ``若 $s$ 是 $S$ 的元素, 则 $s$ 是 $T$ 的元素''; 也就是说, $S$ 的所有元素都是 $T$ 的元素; 也就是说, 正如所许, $S \subseteq T$.

注意, 对所有集合 $S$, 有 $\emptyset \subseteq S$ 且 $S \subseteq S$.

若 $S \subseteq T$ 且 $T \subseteq S$, 则 $S = T$.

符号 $|S|$ 表示 $S$ 的元素个数 (当这个数有限时); 否则记 $|S| = \infty$. 若 $S$ 与 $T$ 有限, 则
\[
\index{Lie@Lie!Sophus@Sophus!daishu@代数(algebra)}	S \subseteq T \implies |S| \leq |T|.
\]
一个集合 $S$ 的全体子集构成一个集合, 称为 $S$ 的幂集%
\index{miji@幂集(power set)}
\index{juzhendexingkongjian@矩阵的行空间(row space of a matrix)}\index{juzhendexingkongjian@矩阵的行空间(row space of a matrix)!huan@环(ring)}\index{mi@幂(power)!ziyu@子域(subfield)!bukeyue@不可约(irreducible)}\index{mi@幂(power)!feilinglixiangzai PID zhongjida@非零理想在 PID 中极大(nonzero ideals are maximal in a PID)!bukeyue@不可约(irreducible)}\index{mi@幂(power)!lixiang@理想(ideal)!bukeyue@不可约(irreducible)}\index{mi@幂(power)!zhenghuanzhongdeyuan@整环中的元(element in an integral domain)!bukeyue@不可约(irreducible)}\index{mi@幂(power)!xiangban@相伴(associated)!mode@模的(of a module)!bukeyue@不可约(irreducible)}\index{mi@幂(power)!jishu@级数(series)!bukeyue@不可约(irreducible)}\index{mi@幂(power)!jishu@级数(series)}\index{mi@幂(power)}(或``部分集合之集''). 例如, 空集 $\emptyset$ 的幂集由一个元素构成: $\{\emptyset\}$. $S$ 的幂集记作 $\mathcal{P}(S)$; 常见的替代记号是 $2^S$, 而且当 $S$ 有限时确有 $|\mathcal{P}(S)| = 2^{|S|}$ (参见习题~\ref{ex:two-to-power}).\index[sym1]{45@$\mathcal{P}(S)$}

\subsection{集合间的运算}
\label{subsec:operations}

手头有了几个集合之后, 就可以通过某些标准运算得到更多的集合. 这里有几个:
\begin{itemize}
	\item $\cup$: 并;
	\item $\cap$: 交;
	\item $\smallsetminus$: 差;
	\item $\amalg$: 无交并%
\index{wujiaobing@无交并(disjoint union)}
;
	\item $\times$: (Cartesian) 积%
\index{jij@积(product)}
;
	\item 以及``模等价关系取商''这一重要概念.
\end{itemize}
这些运算中的大多数读者应当已经熟悉: 例如,
\[
	\{1,2,4\} \cup \{3,4,5\} = \{1,2,3,4,5\},
\]
而
\[
	\{1,2,4\} \smallsetminus \{3,4,5\} = \{1,2\}.
\]
用令人回忆起臭名昭著的``新数学''的维恩图%
\index{weientu@维恩图(Venn diagram)}
表示如下:
\index{zhongxin@中心(center)!huande@环的(of a ring)}\index{zhongxin@中心(center)!qunde@群的(of a group)}\begin{center}
\begin{tikzpicture}[scale=0.72,line width=0.5pt]
\begin{scope}
\draw[pattern=north east lines] (0,0) ellipse (1.5 and .85);
\draw[pattern=north west lines] (1.7,0) ellipse (1.5 and .85);
\node at (-2.1,.7) {$S$}; \node at (3.8,.7) {$T$};
\end{scope}
\foreach \x/\opname in {-5.4/{S\cup T},.85/{S\cap T},7.1/{S\smallsetminus T}} {
\begin{scope}[shift={(\x,-4.2)}]
\draw[pattern=north east lines] (0,0) ellipse (1.5 and .85);
\draw[pattern=north west lines] (1.7,0) ellipse (1.5 and .85);
\node at (.85,-1.5) {$\opname$};
\end{scope}}
\begin{scope}[shift={(-5.4,-4.2)}]
\draw[line width=1.6pt] ({1.5*cos(55.48)},{.85*sin(55.48)}) arc[start angle=55.48,end angle=304.52,x radius=1.5,y radius=.85];
\draw[line width=1.6pt] ({1.7+1.5*cos(235.48)},{.85*sin(235.48)}) arc[start angle=235.48,end angle=484.52,x radius=1.5,y radius=.85];
\end{scope}
\begin{scope}[shift={(.85,-4.2)}]
\draw[line width=1.6pt] ({1.5*cos(-55.48)},{.85*sin(-55.48)}) arc[start angle=-55.48,end angle=55.48,x radius=1.5,y radius=.85];
\draw[line width=1.6pt] ({1.7+1.5*cos(124.52)},{.85*sin(124.52)}) arc[start angle=124.52,end angle=235.48,x radius=1.5,y radius=.85];
\end{scope}
\begin{scope}[shift={(7.1,-4.2)}]
\draw[line width=1.6pt] ({1.5*cos(55.48)},{.85*sin(55.48)}) arc[start angle=55.48,end angle=304.52,x radius=1.5,y radius=.85];
\draw[line width=1.6pt] ({1.7+1.5*cos(235.48)},{.85*sin(235.48)}) arc[start angle=235.48,end angle=124.52,x radius=1.5,y radius=.85];
\end{scope}
\end{tikzpicture}
\end{center}
(其中黑色实线轮廓指明了参与该运算的集合).

其中几种运算可以用逻辑语言透明地写出: 例如,
\[
	s \in S \cap T \iff (s \in S \ \text{ 并且 } \ s \in T).
\]
两个集合 $S$ 与 $T$ 称为不相交%
\index{buxiangjiao@不相交(disjoint)}
, 如果 $S \cap T = \emptyset$, 也就是说, 没有``同时''属于二者的元素.

子集 $T$ 在集合 $S$ 中的补集%
\index{buji@补集(complement)}
是差集 $S \smallsetminus T$, 它由 $S$ 中所有不在 $T$ 中的元素构成. 例如, 偶整数集在 $\Z$ 中的补集就是奇整数集.

运算 $\amalg$, $\times$ 以及模等价关系取商则稍微神秘一些, 仔细琢磨它们是极有教益的. 我们将先以特别朴素的方式考察它们, 待获得了更多语言、能够从更成熟的视角审视它们之后, 再回过头来讨论.

\subsection{无交并、积}
\index{ji@积(product)!mode@模的(of a module)}\index{ji@积(product)!qunde@群的(of a group)}\index{ji@积(product)!zai Abel fanchouzhong@在 Abel 范畴中(in an abelian category)}\index{ji@积(product)!ziyou@自由(free)}\index{ji@积(product)!xianwei@纤维(fiber)}\label{subsec:disjoint-unions-products}

这些运算有一个问题: 它们的输出可能无法定义为一个集合, 而只能定义为一个``相差集合之同构意义上的''集合, 也就是说, 相差双射意义上的集合. 要把这话说清楚, 就必须谈论映射; 我们马上就来做这件事.

粗略地说, 两个集合 $S$ 与 $T$ 的无交并是一个集合 $S \amalg T$, 它这样得到: 先造出 $S$ 与 $T$ 的``副本'' $S'$ 与 $T'$, 满足 $S' \cap T' = \emptyset$, 然后取 $S'$ 与 $T'$ 的 (普通) 并. 细心的读者会感到不安, 因为这份``配方''并未定义出一个集合: 无论``造一个集合的副本''是什么意思, 肯定有许多种造法. 这一歧义将在下文澄清.

尽管如此, 请注意, 即便在这样站不住脚的基础上, 关于 $S \amalg T$ 我们仍能说点什么: 例如, 若 $S$ 由 $3$ 个元素组成, $T$ 由 $4$ 个元素组成, 那么读者应当 (而且能够正确地) 期望 $S \amalg T$ 由 $7$ 个元素组成.

积也被同样类型的歧义所困扰, 但幸运的是有一个方便的约定, 使我们能够写下代表两个集合 $S$ 与 $T$ 之积的``某一个''集合: 给定 $S$ 与 $T$, 我们令 $S \times T$ 为由 $S$ 与 $T$ 的元素构成的有序对%
\index{youxudui@有序对(ordered pair)}
$(s,t)$ 组成的集合%
\footnote{有序对 $(s,t)$ 可以用集合来定义: 令 $(s,t) = \{s, \{s,t\}\}$. 这样既携带了元素 $s$ 与 $t$ 的信息, 又表达了 $s$ 的特殊地位 (= 序对的第一个元素).}%
:
\[
	S \times T := \{(s,t) \text{ 使得 } s \in S,\ t \in T\}.
\]
于是, 若 $S = \{1,2,3\}$, $T = \{3,4\}$, 则
\[
	S \times T = \{(1,3), (1,4), (2,3), (2,4), (3,3), (3,4)\}.
\]
一个更精致的例子: $\R \times \R$ 是实数对的集合, 它 (正如读者在微积分中学到的) 是平面的一种很好的表示. 集合 $\Z \times \Z$ 则可以表示为这个平面中恰好具有整数坐标的那些点. 顺便一提, 人们常把这两个集合记作 $\R^2$, $\Z^2$; 类似地, 集合与其自身的积 $A \times A$ 常记作 $A^2$.

若 $S$ 与 $T$ 是有限集, 则显然 $|S \times T| = |S|\,|T|$.

还应注意, 我们可以用积来按需造出无交并所需的``副本'': 例如, 可以令 $S' = \{0\} \times S$, $T' = \{1\} \times T$, 这就保证了 $S'$ 与 $T'$ 不相交 (为什么?); 并且有一种明显的方式去``认同'' $S$ 与 $S'$、$T$ 与 $T'$. 再说一遍, 把这话说精确还需要一点词汇.

运算 $\cup$, $\cap$, $\amalg$, $\times$ 可以推广到整``族''集合上的运算: 例如, 若 $S_1, \ldots, S_n$ 是集合, 我们写
\[
\bigcap_{i=1}^{n} S_i = S_1 \cap S_2 \cap \dots \cap S_n
\]
表示由那些同时是 $S_1, \ldots, S_n$ 中每一个集合的元素组成的集合; 其他运算类似. 但请注意, 虽然从定义看显然有
\[
S_1 \cup S_2 \cup S_3 = (S_1 \cup S_2) \cup S_3 = S_1 \cup (S_2 \cup S_3),
\]
但
\[
S_1 \times S_2 \times S_3, \quad (S_1 \times S_2) \times S_3, \quad S_1 \times (S_2 \times S_3)
\]
这些集合应当在何种意义上被``认同'', 却不那么清楚 (此时最左边的集合可以仿照两个集合的情形, 定义为由 $S_1$, $S_2$, $S_3$ 中元素构成的``有序三元组''之集). 事实上, 又一次地, 只有在掌握了映射的语言之后, 我们才能真正理解这类陈述. 不过, 正如读者大概期望的那样, 所有这类陈述最终都被证实为真; 由于这一幸运的情况, 我们可以稍稍随便一些, 对这类微妙之处一笔带过.

更一般地, 若 $\mathcal{S}$ 是一个由集合构成的集合, 我们可以考虑
\[
\bigcup_{S \in \mathcal{S}} S, \quad \bigcap_{S \in \mathcal{S}} S, \quad \coprod_{S \in \mathcal{S}} S, \quad \prod_{S \in \mathcal{S}} S,
\]
即 $\mathcal{S}$ 中所有集合的并、交、无交并与积. 这些定义涉及一些重要的微妙之处: 例如, 若所有 $S \in \mathcal{S}$ 都非空, 是否能推出 $\prod_{S \in \mathcal{S}} S$ 非空? 读者大概认为是显然的, 但 (若 $\mathcal{S}$ 无限) 这是相当棘手的问题, 它等价于选择公理%
\index{xuanzegongli@选择公理(axiom of choice)}
.

大体而言, 这些微妙之处并不影响本课程的内容; 待它们与手头的问题变得更为相关时, 我们会 (在 \S V.\ref{sec:zorn-lemma}) 部分地与它们和解%
\footnote{读者在此后的若干习题中将时不时需要动用选择公理, 即便在我们回头讨论这些问题之前也是如此; 但这多半会发生在觉察阈值之下, 而这恰恰是它应有的样子.}%
.

\subsection{等价关系、划分、商}
\index{guanxi@关系(relation)!jie@阶(order)}\index{guanxi@关系(relation)!dengjia@等价(equivalence)}\index{guanxi@关系(relation)!yuqunjiegouxiangrong@与群结构相容(compatible with a group structure)}\label{subsec:equivalence-relations}

直观地说, 一个集合 $S$ 的元素之间的关系%
\index{guanxi@关系(relation)}
, 是 $S$ 的元素的某些选择之间的某种特殊亲缘. 例如, $\Z$ 上的关系 $<$ 是比较两个整数大小的方式: 由于 $2 < 5$, 在这个意义下 $2$``与''$5$ 有关系, 而 $5$ 在同一意义下与 $2$ 没有关系.

就一切实用目的而言, 一个关系``意指''什么, 完全由集合中哪些元素与哪些元素有关系所刻画. 如果我们拥有一份由所有使得 $a < b$ 的整数对 $(a,b)$ 组成的完整清单, 那么关于 $\Z$ 上的 $<$ 我们就无所不知了. 例如, $(2,5)$ 是这样的一对, 而 $(5,2)$ 不是.

这就引出了关系概念的完全直截了当的定义: 集合 $S$ 上的关系就是积 $S \times S$ 的一个子集 $R$. 若 $(a,b) \in R$, 我们说 $a$ 与 $b$ ``有关系 $R$'', 并写作
\[
a \mathrel{R} b.
\]
人们常为关系使用更花哨的符号, 诸如 $<$, $\leq$, $=$, $\sim$, \ldots.

性质良好的关系的原型是``='', 它对应于``对角线''
\[
\{(a,b) \in S \times S \mid a = b\} = \{(a,a) \mid a \in S\} \subseteq S \times S.
\]
这个非常特殊的关系的三条性质被证明是特别重要的: 若暂时用 $\sim$ 记相等关系 $=$, 则 $\sim$ 满足
\begin{itemize}
	\item 自反性: $(\forall a \in S)\ a \sim a$;
	\item 对称性: $(\forall a \in S) (\forall b \in S)\ a \sim b \implies b \sim a$;
	\item 传递性: $(\forall a \in S) (\forall b \in S) (\forall c \in S),\ (a \sim b \text{ 并且 } b \sim c) \implies a \sim c$.
\end{itemize}
也就是说, 每个 $a$ 都等于它自身; 若 $a$ 等于 $b$, 则 $b$ 等于 $a$; 等等.

\begin{definition}[等价关系]
\label{def:equiv-rel}
集合 $S$ 上的等价关系%
\index{dengjiaguanxi@等价关系(equivalence relation)}
是指满足上述三条性质的任意关系 $\sim$.
\end{definition}

用 $S \times S$ 中相应的子集 $R$ 的语言来说, ``自反性''说的是对角线包含于 $R$; ``对称性''说的是 $R$ 关于对角线翻转不变 (即每个 $(a,b)$ 换成 $(b,a)$); 而遗憾的是, ``传递性''没有同样漂亮的图示翻译.

$S$ 上的等价关系这一数据, 结果等价于一类乍看之下稍有不同的信息, 即 $S$ 的一个划分%
\index{huafen@划分(partition)}
. $S$ 的划分是 $S$ 的一族两两不交的非空子集, 其并为 $S$: 例如,
\[
\mathscr{P} = \{\{1,4,7\}, \{2,5,8\}, \{3,6\}, \{9\}\}
\]
是集合
\[
\{1,2,3,4,5,6,7,8,9\}
\]
的一个划分.

以下是如何从 $S$ 上的关系 $\sim$ 得到 $S$ 的一个划分: 对每个元素 $a \in S$, $a$ 的 (相对于 $\sim$ 的) 等价类%
\index{dengjialei@等价类(equivalence class)}
是 $S$ 的子集
\[
[a]_{\sim} := \{b \in S \mid b \sim a\};
\]
于是这些等价类构成 $S$ 的一个划分 $\mathscr{P}_{\sim}$ (习题~\ref{ex:partition-from-relation}).

反之 (习题~\ref{ex:relation-from-partition}), 每个划分 $\mathscr{P}$ 都以这种方式对应于某个等价关系. 因此, ``$S$ 上的等价关系''与``$S$ 的划分''这两个概念实质上是等价的.

现在我们可以把 $\mathscr{P}_{\sim}$ 视为一个集合 (其元素是相对于 $\sim$ 的等价类). 这就是 \S\ref{subsec:operations} 中提到的取商运算.

\begin{definition}[商]
\label{def:quotient}
集合 $S$ 关于等价关系 $\sim$ 的商%
\index{shangji@商(quotient)}
是 $S$ 的元素相对于 $\sim$ 的等价类之集
\[
S / {\sim} := \mathscr{P}_{\sim}.
\]
\end{definition}

\begin{example}
\label{ex:mod-two}
取 $S = \Z$, 并令 $\sim$ 为由
\[
a \sim b \iff a - b \text{ 是偶数}
\]
定义的关系. 则 $\Z / {\sim}$ 由两个等价类构成:
\[
\Z / {\sim} = \{[0]_{\sim}, [1]_{\sim}\}.
\]
的确, 每个整数 $b$ 要么是偶数 (从而 $b - 0$ 是偶数, 故 $b \sim 0$, 即 $b \in [0]_{\sim}$), 要么是奇数 (从而 $b - 1$ 是偶数, 故 $b \sim 1$, 即 $b \in [1]_{\sim}$). 这当然是模算术的起点, 稍后我们将对它作应尽的详细介绍 (\S II.\ref{subsec:cyclic-groups}).

思考这一运算的一种方式是: 等价关系在商中``变成相等'': 也就是说, 商 $S/{\sim}$ 中的两个元素相等, 当且仅当 $S$ 中相应的元素有关系 $\sim$. 换言之, 取商是一种把任意等价关系变成相等关系的办法. 这一观察稍后将在``范畴的语言''中进一步形式化 (\S\ref{subsec:quotients-up}).
\end{example}

\subsection*{习题}
\phantomsection
\addcontentsline{toc}{subsection}{习题}

标有 $\triangleright$ 的习题被正文引用; 标有 $\neg$ 的习题被其他习题引用. 引用它们的习题与节列在该习题后面的方括号中; 如有必要, 请参阅前言中的进一步说明.

\begin{prob}
\label{ex:russell}
查找关于罗素悖论的讨论, 并把它弄懂.
\end{prob}

\begin{prob}
\label{ex:partition-from-relation}
$\triangleright$ 证明: 若 $\sim$ 是集合 $S$ 上的一个等价关系, 则 \S\ref{subsec:equivalence-relations} 中定义的相应族 $\mathscr{P}_{\sim}$ 确实是 $S$ 的一个划分: 也就是说, 其元素非空、两两不交, 且其并为 $S$. [\S\ref{subsec:equivalence-relations}]
\end{prob}

\begin{prob}
\label{ex:relation-from-partition}
$\triangleright$ 给定集合 $S$ 上的一个划分 $\mathscr{P}$, 说明如何在 $S$ 上定义一个等价关系\erratumfootnote{根据官方勘误, 此处应为``等价关系''(equivalence relation).} $\sim$, 使得 $\mathscr{P}$ 正是与之相应的划分. [\S\ref{subsec:equivalence-relations}]
\end{prob}

\begin{prob}
\label{ex:equiv-rel-123}
在集合 $\{1,2,3\}$ 上可以定义多少个不同的等价关系?
\end{prob}

\begin{prob}
\label{ex:refl-symm-not-trans}
给出一个自反且对称但不传递的关系的例子. 如果试图用这个关系在集合上定义划分, 会发生什么? (提示: 思考第二个问题将有助于回答第一个问题.)
\end{prob}

\begin{prob}
\label{ex:quotient-R}
$\triangleright$ 在实数集 $\R$ 上定义关系 $\sim$: 规定 $a \sim b \iff b - a \in \Z$. 证明这是一个等价关系, 并对 $\R / {\sim}$ 找出一个``令人信服''的描述. 对平面 $\R \times \R$ 上通过规定 $(a_1, a_2) \approx (b_1, b_2) \iff b_1 - a_1 \in \Z$ 且 $b_2 - a_2 \in \Z$ 定义的关系 $\approx$ 做同样的事情. [\S~II.\ref{subsec:canon-quotient}, II.\ref{ex:torus}]
\end{prob}

\section{集合间的映射}
\label{sec:functions}

\subsection{定义}
\label{subsec:functions-def}

对于本书引入的几乎每一种结构, 我们都将遵循一条共同的主线: 既设法理解这种结构本身, 也理解这种结构的不同实例之间可能如何相互作用.

集合通过映射%
\index{yingshe@映射(map)}
相互作用. 人们难免想用``动态''的措辞来思考从集合 $A$ 到集合 $B$ 的映射 $f$, 把它当作``从 $A$ 走到 $B$''的一种方式. 与关系的情形类似, 把这个概念形式化而无需求助任何给定 $f$ 的深层``含义''是直截了当的: 关于映射 $f$ 可以知道的一切, 都被如下信息所刻画, 即 $A$ 中任一给定元素 $a$ 的像是 $B$ 中哪个元素 $b$. 这个信息不过是 $A \times B$ 的一个子集:
\[
\Gamma_f := \{(a,b) \in A \times B \mid b = f(a)\} \subseteq A \times B.
\]
集合 $\Gamma_f$ 称为 $f$ 的图像%
\index{tuxiang@图像(picture)}
; 严格地说, 映射``就是''它的图像%
\footnote{更精确地说, 是图像 $\Gamma_f$ 连同 $f$ 的源 $A$ 与目标 $B$ 的信息; 这些都属于映射的数据的一部分.}%
.

并非 $A \times B$ 的所有子集 $\Gamma \subseteq A \times B$ 都对应于 (即``是'') 映射: 我们需要对映射的图像施加一条要求, 它可以表述如下:
\[
(\forall a \in A) (\exists! b \in B)\ (a,b) \in \Gamma_f,
\]
或者 (``用函数记号'')
\[
(\forall a \in A) (\exists! b \in B)\ f(a) = b.
\]
也就是说, 映射必须把 $A$ 的每个元素 $a$ 送到 $B$ 的恰好一个元素, 而且该元素依赖于 $a$. 诸如 $\pm\sqrt{x}$ 这样的``多值函数'' (它们在, 比如, Riemann 面的研究中非常重要) 在这个意义下不是映射.

要宣布 $f$ 是从集合 $A$ 到集合 $B$ 的映射, 就写 $f : A \to B$, 或画出下图 (``图''):
\[
A \xrightarrow{f} B.
\]
映射 $f : A \to B$ 在元素 $a \in A$ 上的作用有时用一个``带装饰的''箭头表示, 如
\[
a \mapsto f(a).
\]
从一个集合 $A$ 到一个集合 $B$ 的所有映射的汇集本身也是一个集合%
\footnote{这是另一种``集合间的运算'', 未列在 \S\ref{subsec:operations} 中. 你能看出我们为什么把这个集合记作 $B^A$ 吗? (参见习题~\ref{ex:card-ba}.)}\index[sym1]{06@$B^A$}%
, 记作 $B^A$. 如果我们认真对待``映射其实与它的图像是同一回事''这一观念, 那么可以把 $B^A$ 视为 $A \times B$ 的幂集的一个 (特殊) 子集.

每个集合 $A$ 都自带一个非常特殊的映射, 其图像是 $A \times A$ 中的对角线: $A$ 上的恒同映射%
\index{hengtong@恒同(identity)}
\[
\id_A : A \to A
\]
由 $(\forall a \in A)\ \id_A(a) = a$ 定义. 更一般地, 集合 $A$ 的任意子集 $S$ 的包含关系确定一个映射 $S \to A$, 它把 $S$ 的每个元素 $s$ 送到 $A$ 中``它自己''.

若 $S$ 是 $A$ 的子集, 我们用 $f(S)$ 表示 $B$ 的子集
\[
f(S) := \{b \in B \mid (\exists a \in S)\ b = f(a)\}.
\]
也就是说, $f(S)$ 是 $B$ 的由 $S$ 的元素在映射 $f$ 下的全部像组成的子集. 其中最大者, 即 $f(A)$, 称为 $f$ 的像%
\index{xiang@像(image)}
, 记作``$\Image f$''.

另外, $f|_S$ 表示 $f$ 在子集 $S$ 上的``限制''%
\index{xianzhi@限制(restriction)}
: 这是映射 $S \to B$, 定义为
\[
(\forall s \in S): \quad f|_S(s) = f(s).
\]
也就是说, $f|_S$ 是复合 (按下一小节解释的意义) $f \circ i$, 其中 $i : S \to A$ 是包含映射. 注意 $f(S) = \Image (f|_S)$.

\subsection{例子: 多重集、加标集合}
\label{subsec:multisets-indexed}

\S\ref{subsec:sets} 中简略提到的``多重集''是借助映射容易形式化的概念的一个简单例子. 多重集可以这样定义: 给定一个从 (普通) 集合 $A$ 到正%
\index{xianxingbianhuandetezhengzhi@线性变换的特征值(eigenvalue of a linear transformation)}\index{xianxingbianhuandetezhengzhi@线性变换的特征值(eigenvalue of a linear transformation)!dechongshu@的重数(multiplicity of)}\footnote{一些文献允许 $0$ 作为可能的重数.}%
整数集 $\N^*$ 的映射; 若 $m : A \to \N^*$ 是这样一个映射, 相应的多重集由元素 $a \in A$ 组成, 每个元素取 $m(a)$ 次. 于是, 多重集 $\{a,a,a,b,b,b,b,b,c\}$ 实质上就是满足 $m(a) = 3$, $m(b) = 5$, $m(c) = 1$ 的映射 $m : \{a,b,c\} \to \N^*$. 与普通集合一样, 元素列出的次序不属于多重集所携带的信息. 包含、并等简单的集合论概念可以直截了当地推广到多重集. 关于多重集的另一种观点, 见习题~\ref{ex:mset-category}.

另一个例子是``指标''的使用. 如果我们写``令 $a_1, \ldots, a_n$ 为整数\ldots'', 我们真正的意思是``考虑一个映射 $a : \{1, \ldots, n\} \to \Z$ \ldots'', 并理解 $a_i$ 是值 $a(i)$ 的简写 (对 $i = 1, \ldots, n$). 人们难免把加标集合%
\index{jiabiaojih@加标集合(indexed set)}
$\{a_i\}_{i \in I}$ 简单地想成一个集合, 其元素恰好被记作 $a_i$, 而 $i$ 取遍某个``指标集'' $I$; 但这样的加标集合更严格地说是一个映射 $I \to A$, 其中 $A$ 是某个我们从中提取元素 $a_i$ 的集合. 例如, 这使得我们可以把 $a_0$ 与 $a_1$ 视为 $\{a_i\}_{i \in \N}$ 的两个不同元素, 即便碰巧作为目标集合 $A$ 的元素有 $a_0 = a_1$.

\index{xiangliang@向量(vector)!cong@丛(bundle)}这类微妙之处很容易被忽视, 而且记号上的某些滥用是常见的、通常也是无害的. 这些区分在 (例如) 关于向量组%
\index{xiangliangzu@向量组(vectors)}
\index{xianxing（wu）guan@线性（无）关(linearly dependent/independent)}\index{xianxing@线性(linear)}线性无关%
\index{xianxingwuguan@线性无关(linearly independent)}
的讨论中起作用; 参见 \S VI.\ref{subsec:linear-independence-bases}.

\subsection{映射的复合}
\label{subsec:composition}

映射可以复合: 若 $f : A \to B$ 与 $g : B \to C$ 是映射, 则由
\begin{equation*}
(\forall a \in A) \quad (g \circ f)(a) := g(f(a)) \tag{$*$}
\end{equation*}
定义的运算 $g \circ f$ 也是映射. 也就是说, 我们先利用 $f$ 从 $A$ 到 $B$, 再应用 $g$ 到达 $C$. 用图形可以画出诸如
\[
\begin{tikzcd}[column sep=large]
A \ar[r,"f"] \ar[rr,bend right=25,"g\circ f"'] & B \ar[r,"g"] & C
\end{tikzcd}
\qquad\text{或}\qquad
\begin{tikzcd}
& B \ar[dr,"g"] & \\
A \ar[ur,"f"] \ar[rr,"g\circ f"'] & & C
\end{tikzcd}
\]
的图. 用这种图形表示 (比如) 由映射相连的集合的汇集, 称为图%
\index{tu@图(graph)}
. 我们将画许多图, 而且是在比眼下这个语境一般得多的语境中.

我们说上面画的图``交换'', 或``是交换的''%
\index{jiaohuantu@交换图(commutative diagram)}
, 意思是: 若我们从 $A$ 出发, 沿图中规定的两种可能路径之一走到 $C$, 沿途应用所遇到的映射, 得到的结果相同. 这恰恰就是陈述 $(*)$ 的内容.

复合是结合的: 也就是说, 若 $f : A \to B$, $g : B \to C$, $h : C \to D$ 是映射, 则 $h \circ (g \circ f) = (h \circ g) \circ f$. 用图形表示, 图
\begin{center}
\begin{tikzcd}[column sep=1.6em]
	A \ar[r, "f"] \ar[rr, bend right=22, "g \circ f"'] & B \ar[r, "g"] \ar[rr, bend left=22, "h \circ g"] & C \ar[r, "h"] & D
\end{tikzcd}
\end{center}
交换. 从复合的定义看, 这个重要的观察应当是完全显然的.

恒同映射关于复合是非常特殊的: 若 $f : A \to B$ 是任意映射, 则 $\id_B \circ f = f$ 且 $f \circ \id_A = f$. 用图形表示, 图
\[
\begin{tikzcd}[column sep=large]
A \ar[r,"f"] \ar[rr,bend right=25,"f"'] & B \ar[r,"\id_B"] & B
\end{tikzcd}
\qquad
\begin{tikzcd}[column sep=large]
A \ar[r,"\id_A"] \ar[rr,bend right=25,"f"'] & A \ar[r,"f"] & B
\end{tikzcd}
\]
交换.

\subsection{单射、满射、双射}
\label{subsec:inj-surj-bij}

有几类特殊的映射值得突出标出:
\begin{itemize}
	\item 映射 $f : A \to B$ 称为单的%
(或单射%
\index{danshe@单射(injective)}
, 或一对一), 如果
	\[
\index{Lie@Lie!Sophus@Sophus!kuohao@括号(bracket)}	a' \neq a'' \implies f(a') \neq f(a''):
	\]
	也就是说, $f$ 把不同的元素送到不同的元素%
	\footnote{人们常用逆否 (因而等价) 的形式来检验这个定义, 即 $(\forall a' \in A)(\forall a'' \in A)\ f(a') = f(a'') \implies a' = a''$.}%
	.

	\item 映射 $f : A \to B$ 称为满的%
\index{mande@满的(full)}
(或满射%
\index{manshe@满射(surjective)}
, 或映上), 如果
	\[
	(\forall b \in B) (\exists a \in A) \quad b = f(a):
	\]
	也就是说, $f$``覆盖整个 $B$''; 更精确地说, $\Image f = B$.
\end{itemize}
单射常画作 $\hookrightarrow$; 满射常画作 $\twoheadrightarrow$.

若 $f$ 既是单的又是满的, 就称它是双射%
\index{shuangshe@双射(bijective)}
(或一一对应, 或集合之同构%
\index{tonggou@同构(isomorphism)}
). 这时我们常写 $f : A \xrightarrow{\sim} B$, 或
\[
A \cong B,
\]
并说 $A$ 与 $B$ 是``同构''的集合.

当然, 恒同映射 $\id_A : A \to A$ 是双射.

若 $A \cong B$, 即存在双射 $f : A \to B$, 则集合 $A$ 与 $B$ 可以通过 $f$ 被``认同'', 意思是我们能把 $A$ 的元素 $a$ 与 $B$ 中相应的元素 $f(a)$ 精确地匹配起来. 例如, 若 $A$ 是有限集且 $A \cong B$, 则 $B$ 必然也是有限集, 且 $|A| = |B|$.

有了这套术语, 我们就能更好地理解 \S\ref{subsec:operations} 中关于``无交并''的讨论: 所给集合 $A$, $B$ 的``副本'' $A'$, $B'$ 就应当分别是与 $A$, $B$ 同构的集合. \S\ref{subsec:disjoint-unions-products} 末尾提出的制造这种不交``副本''的方案是可行的, 因为 (比如) 由
\[
(\forall a \in A) \quad f(a) = (0, a)
\]
定义的
\[
f : A \to \{0\} \times A
\]
显然是双射.

\subsection{单射、满射、双射: 第二种观点}
\label{subsec:second-viewpoint}

思考这些概念还有另一种富有启发性的方式.

若 $f : A \to B$ 是双射, 则可以``把它的图像翻转'', 定义一个映射
\[
g : B \to A:
\]
也就是说, 恰当 $b = f(a)$ 时规定 $a = g(b)$. ($f$ 既单又满这一事实保证了 $\Gamma_f$ 的翻转按 \S\ref{subsec:functions-def} 给出的定义是某个映射的图像. 请自行验证!)

这个映射 $g$ 有一个非常有趣的性质: 用图形表示,
\[
\begin{tikzcd}[column sep=large]
A \ar[r,"f"] \ar[rr,bend right=25,"\id_A"'] & B \ar[r,"g"] & A
\end{tikzcd}
\qquad
\begin{tikzcd}[column sep=large]
B \ar[r,"g"] \ar[rr,bend right=25,"\id_B"'] & A \ar[r,"f"] & B
\end{tikzcd}
\]
交换; 也就是说, $g \circ f = \id_A$ 且 $f \circ g = \id_B$. 第一个等式告诉我们 $g$ 是 $f$ 的``左逆''%
\footnote{别管在图中 $g$ 被画在 $f$ 的右边这件事---我们说 $g$ 是 $f$ 的左逆, 是因为它写在 $f$ 的左边: $g \circ f = \id_A$.}%
; 第二个等式告诉我们 $g$ 是 $f$ 的``右逆''. 我们干脆说它是 $f$ 的逆%
\index{nid@逆(inverse)}
, 记作 $f^{-1}$. 于是, ``双射有逆''.

\index{ni@逆(inverse)!shiweiyide@是唯一的(is unique)}那逆命题呢? 若一个映射有逆, 它是双射吗? 这是真的, 而且事实上我们能说得精确得多.

\begin{proposition}
\label{prop:left-right-inverse}
设 $A \neq \emptyset$, 并设 $f : A \to B$ 是映射. 则
\begin{enumerate}
	\item $f$ 有左逆当且仅当它是单的;
	\item $f$ 有右逆当且仅当它是满的.
\end{enumerate}
\end{proposition}

\begin{proof}
来证 (1).

($\implies$) 若 $f : A \to B$ 有左逆, 则存在 $g : B \to A$ 使 $g \circ f = \id_A$. 现设 $a' \neq a''$ 是 $A$ 中任意两个不同的元素; 那么
\[
g(f(a')) = \id_A(a') = a' \neq a'' = \id_A(a'') = g(f(a''));
\]
也就是说, $g$ 把 $f(a')$ 与 $f(a'')$ 送到不同的元素. 这迫使 $f(a')$ 与 $f(a'')$ 不同, 表明 $f$ 是单的.

($\Leftarrow$) 现设 $f : A \to B$ 是单的. 为了构造映射 $g : B \to A$, 我们必须对每个元素 $b \in B$ 指定唯一的值 $g(b) \in A$. 为此, 任取一个固定元素 $s \in A$ (由于 $A \neq \emptyset$, 这做得到); 然后令
\[
g(b) := \begin{cases}
	a & \text{若 } b = f(a) \text{ 对某个 } a \in A, \\
\index{xiang@像(image)!tonggouyuyuxiang@同构于余像(isomorphic to coimage)}	s & \text{若 } b \notin \Image f.
\end{cases}
\]
用文字说, 若 $b$ 是 $A$ 中某个元素 $a$ 的像, 就把它送回 $a$; 否则, 把它送到你那个固定的元素 $s$.

上面这个指定确实定义了一个映射, 恰恰因为 $f$ 是单的: 的确, 这保证了每个作为 $A$ 中某个 $a$ 在 $f$ 下的像的 $b$, 都是唯一的 $a$ 的像 ($A$ 中两个不同元素不可能被 $f$ 同时送到 $b$, 因为 $f$ 是单的). 于是每个 $b \in B$ 都被送到 $A$ 中唯一的良定义的元素, 正如对映射所要求的那样.

最后, 映射 $g : B \to A$ 是 $f$ 的左逆. 的确, 若 $a \in A$, 则 $b = f(a)$ 属于第一种类型, 于是被 $g$ 送回 $a$; 也就是说, 对所有 $a \in A$ 有 $g \circ f(a) = a = \id_A(a)$, 正如所需.

(2) 的证明留作习题 (习题~\ref{ex:right-inverse}).
\end{proof}

\begin{corollary}
\label{cor:bijection-inverse}
映射 $f : A \to B$ 是双射当且仅当它有 (双侧的) 逆.
\end{corollary}

这并非完全显然: 若 $f$ 既有左逆又有右逆, 为什么它必有一个同时作为左逆与右逆的逆? 请读者现在尝试自行证明. 我们很快会在 \S\ref{sec:morphisms} 回到这个问题.

\index{lunzheng@论证(argument)}若一个映射是单的但不满, 那么它不会有右逆, 而且, 若源至少含两个元素, 它必然有不止一个左逆 (从命题~\ref{prop:left-right-inverse} 的证明所给的论证来看, 这应当是清楚的). 类似地, 一个满射映射通常会有许多右逆; 它们常被称为截面%
\index{jiermian@截面(section)}
.

命题~\ref{prop:left-right-inverse} 暗示这里有一些深刻的东西. \S\ref{subsec:inj-surj-bij} 中给出的单映射与满映射的定义, 关键依赖于直接处理集合的元素; 而命题~\ref{prop:left-right-inverse} 表明, 事实上这些性质可以通过映射在集合之间``组织''的方式被探测到. 即便我们不知道``元素''是什么意思, 我们仍然可以只借助映射的性质来理解单性、满性 (从而还有集合之同构) 的概念.

\index{fanchou@范畴(category)!Ab@Ab!Fld@Fld}这是一个更``成熟''的观点, 也是当我们谈论范畴时将要力荐的观点. 在某种程度上, 它应当能治愈读者的不适---正如我们在对集合的非正式介绍中所做的那样, 我们谈论了``元素'', 却始终没有定义这些神秘实体究竟是什么.

双射 $f$ 的逆的标准记号是 $f^{-1}$. 这个符号也用于非双射的映射, 但是在稍微不同的语境中: 若 $f : A \to B$ 是任意映射, $T \subseteq B$ 是 $B$ 的子集, 则 $f^{-1}(T)$ 表示 $A$ 的``映到 $T$ 的所有元素''构成的子集; 也就是说,
\[
f^{-1}(T) = \{a \in A \mid f(a) \in T\}.
\]
若 $T = \{q\}$ 由 $B$ 的单个元素构成, 则 $f^{-1}(T)$ (简记为 $f^{-1}(q)$) 称为 $f$ 在 $q$ 上的纤维%
\index{xianwei@纤维(fiber)}
. 于是, 映射 $f : A \to B$ 是双射, 当且仅当它在 $B$ 的所有元素上的纤维都非空 (即 $f$ 是满的), 并且这些纤维实际上是单元素集 (即 $f$ 是单的). 在这种情形下, 记号 $f^{-1}$ 与前面提到的``逆''的记号很好地吻合.

\subsection{单态射与满态射}
\label{subsec:mono-epi}

\index{ji@基(basis)!youxu@有序(ordered)}表达单性与满性还有又一种方式, 它初看比我们迄今所见的更复杂, 但实际上更为基本.

映射 $f : A \to B$ 称为单态射%
\index{dantaishe@单态射(monomorphism)}
(或称单的%
\index{dande@单的(simple)}
), 如果下述成立: 对所有集合 $Z$ 与所有映射 $\alpha', \alpha'' : Z \to A$,
\[
f \circ \alpha' = f \circ \alpha'' \implies \alpha' = \alpha''.
\]

\begin{proposition}
\index{neishe@内射(injective)!modexiaojie@模的消解(resolution of a module)}\index{neishe@内射(injective)!xiaojie@消解(resolution)}\index{neishe@内射(injective)!Abel fanchoudeduixiang@Abel 范畴的对象(object of an abelian category)}\index{neishe@内射(injective)!mo@模(module)}\index{neishe@内射(injective)!hanshu@函数(function)}\index{neishe@内射(injective)!zugouneishe@足够内射(enough injectives)}\index{neishe@内射(injective)!weishu@维数(dimension)}\index{neishe@内射(injective)!Abel qun@Abel 群(abelian group)}\index{neishe@内射(injective)!R-Mod youzugouneishe@R-Mod 有足够内射(R-Mod has enough injectives)}\index{neishe@内射(injective)}\label{prop:mono-injective}
一个映射是单射当且仅当它是单态射.
\end{proposition}

\begin{proof}
若映射的源 $A$ 为 $\emptyset$, 则结论显然成立 (为什么?), 故设 $A \neq \emptyset$\erratumfootnote{根据官方勘误, 此处补充对源集为空集情形的说明并假设 $A \neq \emptyset$.}.

($\implies$) 由命题~\ref{prop:left-right-inverse}, 若映射 $f : A \to B$ 是单射, 则它有左逆 $g : B \to A$. 现设 $\alpha', \alpha''$ 是从另一个集合 $Z$ 到 $A$ 的任意映射, 且
\[
f \circ \alpha' = f \circ \alpha'';
\]
在左边复合 $g$, 并利用复合的结合性:
\[
(g \circ f) \circ \alpha' = g \circ (f \circ \alpha') = g \circ (f \circ \alpha'') = (g \circ f) \circ \alpha'';
\]
由于 $g$ 是 $f$ 的左逆, 这就是说
\[
\id_A \circ \alpha' = \id_A \circ \alpha'',
\]
从而
\[
\alpha' = \alpha'',
\]
这正是断言 $f$ 是单态射所需要的结论.

($\Leftarrow$) 现设 $f$ 是单态射. 这说的是关于任意集合 $Z$ 与任意映射 $Z \to A$ 的某件事; 我们将只动用这个信息中微乎其微的一部分, 取 $Z$ 为任意单元素集 $\{p\}$. 这时指定映射 $\alpha', \alpha'' : Z \to A$ 相当于选取把 $Z$ 的唯一元素 $p$ 送到哪些元素 $a' = \alpha'(p)$, $a'' = \alpha''(p)$. 对 $Z$ 的这一特殊选取, 定义单态射的性质
\[
f \circ \alpha' = f \circ \alpha'' \implies \alpha' = \alpha''
\]
变成
\[
f \circ \alpha'(p) = f \circ \alpha''(p) \implies \alpha' = \alpha'';
\]
也就是说
\[
f(a') = f(a'') \implies \alpha' = \alpha'';
\]
现在, 从 $Z = \{p\}$ 到 $A$ 的两个映射相等当且仅当它们把 $p$ 送到同一个元素, 所以这说的是
\[
f(a') = f(a'') \implies a' = a''.
\]
这必须对所有的 $\alpha'$, $\alpha''$ 成立, 也就是说, 对 $A$ 中所有不同的 $a', a''$ 的选取成立. 换言之, $f$ 必须是单射, 正如所要证的.
\end{proof}

读者现在应当会期望: 存在一个按单态射的风格给出的定义, 它将被证明等价于``满''. 情况确实如此: 这个概念称为满态射%
\index{mantaishe@满态射(epimorphism)}
. 找到它, 并证明它与通常的``满''的定义等价, 留给读者%
\footnote{这是一道特别重要的习题, 我们建议读者把所有繁琐的细节都仔细写出来.}%
\index{mantongtai@满同态(epimorphism)!fenlie@分裂(split)}\index{mantongtai@满同态(epimorphism)!jihede@集合的(of sets)}\index{mantongtai@满同态(epimorphism)!huande@环的(of a ring)}\index{mantongtai@满同态(epimorphism)!mode@模的(of a module)}\index{mantongtai@满同态(epimorphism)!qunde@群的(of a group)}\index{mantongtai@满同态(epimorphism)!zaijiaxing/Abel fanchouzhong@在加性/Abel 范畴中(in an additive/abelian category)}\index{mantongtai@满同态(epimorphism)} (习题~\ref{ex:epimorphism}).

\subsection{基本例子}
\label{subsec:basic-examples}

\index{ji@基(basis)!biaozhun@标准(standard)}集合上的基本运算为我们提供了若干单映射与满映射的重要例子.

\begin{example}[投射]
\label{ex:projections}
\index{zirantoushe@自然投射(natural projection)}设 $A$, $B$ 为非空集合\erratumfootnote{根据官方勘误, 补充 $A, B$ 非空的假设.}. 则有自然投射 $\pi_A$, $\pi_B$:
\begin{center}
\begin{tikzcd}[row sep=1.6em]
	& A \times B \ar[ld, two heads, "\pi_A"'] \ar[rd, two heads, "\pi_B"] & \\
	A & & B
\end{tikzcd}
\end{center}
由
\[
\pi_A((a,b)) := a, \qquad \pi_B((a,b)) := b
\]
对一切 $(a,b) \in A \times B$ 定义. 这两个映射都 (显然) 是满的.
\end{example}

\begin{example}[含入无交并]
\label{ex:inclusions}
类似地, 存在从 $A$ 与 $B$ 到无交并的自然单射:
\begin{center}
\begin{tikzcd}[row sep=1.4em]
	A \ar[rd, hook] & & B \ar[ld, hook'] \\
	& A \amalg B &
\end{tikzcd}
\end{center}
它把 $a \in A$ (相应地, $b \in B$) 送到 $A \amalg B$ 中 $A$ 的同构副本 $A'$ (相应地, $B$ 的同构副本 $B'$) 中相应的元素.
\end{example}

\begin{example}[典范投影]
\label{ex:canonical-projection}
若 $\sim$ 是集合 $A$ 上的等价关系, 则有 (显然为满的) 典范投影%
\index{dianfantoushe@典范投影(canonical projection)}
\[
\begin{tikzcd}
A \ar[r,two heads] & A/{\sim}
\end{tikzcd}
\]
它把每个 $a \in A$ 送到其等价类 $[a]_{\sim}$.
\end{example}

\subsection{典范分解}
\label{subsec:canonical-decomposition}

我们把注意力集中在单映射与满映射上的理由是: 它们提供了构造出任何映射的基本``砖块''.

为了看到这一点, 我们观察到: 每个映射 $f : A \to B$ 按如下方式确定 $A$ 上的一个等价关系 $\sim$: 对一切 $a', a'' \in A$,
\[
a' \sim a'' \iff f(a') = f(a'').
\]
(读者应当验证这确实是一个等价关系.)

\begin{theorem}
\label{thm:canonical-decomposition}
设 $f : A \to B$ 是任意映射, 并如上定义 $\sim$. 则 $f$ 有如下分解:
\[\begin{tikzcd}
	A & {(A/{\sim})} & {\Image f} & B
	\arrow[two heads, from=1-1, to=1-2]
	\arrow["f", curve={height=-12pt}, from=1-1, to=1-4]
	\arrow["\sim", from=1-2, to=1-3]
	\arrow["\widetilde f"', draw=none, from=1-2, to=1-3]
	\arrow[hook, from=1-3, to=1-4]
\end{tikzcd}\]
其中第一个映射是典范投影 $A \to A/{\sim}$ (如例~\ref{ex:canonical-projection} 中那样), 第三个映射是包含 $\Image f \subseteq B$, 而中间的双射 $\tilde{f}$ 由
\[
\tilde{f}([a]_{\sim}) := f(a)
\]
对一切 $a \in A$ 定义.
\end{theorem}

定义 $\tilde{f}$ 的公式立即使图交换可见; 因此, 为了证明这个定理, 我们只须验证:
\begin{itemize}
	\item 该公式确实定义了一个映射;
	\item 该映射事实上是双射.
\end{itemize}
第一项属于一类极其重要的验证. 所给的 $\tilde{f}$ 的公式带有一个巨大的内在歧义: $A/{\sim}$ 中的同一个元素可能是 $A$ 中许多元素的等价类; 应用 $\tilde{f}$ 的公式需要从中选取一个元素并对它应用 $f$. 我们必须证明这一操作的结果与该选取无关: 也就是说, 为该等价类选取代表元的一切可能方式都导致同一个结果.

我们把这类情形概括为: 必须验证 $\tilde{f}$ 是良定义的%
\index{liangdingyi@良定义(well-defined)}
. 今后我们会经常需要检验所考虑的运算是否良定义, 其语境与此处这个典型例子十分相似.

\begin{proof}
把上文讨论的第一项展开来, 我们必须验证: 对一切 $a', a'' \in A$,
\[
[a']_{\sim} = [a'']_{\sim} \implies f(a') = f(a'').
\]
而 $[a']_{\sim} = [a'']_{\sim}$ 意味着 $a' \sim a''$, 并且 $\sim$ 的定义当初就是这样精心设计的: 使得这句话恰好意味着这里所需要的 $f(a') = f(a'')$. 所以 $\tilde{f}$ 确实是良定义的.

为验证第二项, 即 $\tilde{f} : A/{\sim} \to \Image f$ 是双射, 我们明确地检验 $\tilde{f}$ 是单的且是满的.

单性: 若 $\tilde{f}([a']_{\sim}) = \tilde{f}([a'']_{\sim})$, 则由 $\tilde{f}$ 的定义有 $f(a') = f(a'')$; 于是由 $\sim$ 的定义有 $a' \sim a''$, 从而 $[a']_{\sim} = [a'']_{\sim}$. 因此
\[
\tilde{f}([a']_{\sim}) = \tilde{f}([a'']_{\sim}) \implies [a']_{\sim} = [a'']_{\sim},
\]
单性得证.

满性: 给定任意 $b \in \Image f$, 存在元素 $a \in A$ 使 $f(a) = b$. 于是由 $\tilde{f}$ 的定义有
\[
\tilde{f}([a]_{\sim}) = f(a) = b.
\]
由于 $b$ 是 $\Image f$ 中的任意元素, 这表明 $\tilde{f}$ 是满的, 正如所需.
\end{proof}

定理~\ref{thm:canonical-decomposition} 表明: 每个映射都是一个满射、接一个同构、再接一个单射的复合. 这个结果的证明是平凡的, 但它具有相当的重要性, 因为它是本书中将多次出现的一类情形的原型. 它还会不时地重新出现, 名字诸如``第一同构定理''.

\subsection{澄清}
\label{subsec:clarification}

最后, 我们可以开始澄清 \S\ref{subsec:disjoint-unions-products} 中关于无交并、积与商的一条评注. 我们对 $A \amalg B$ 的定义是两个不交集 $A'$, $B'$ 的 (约定的) 并, 其中 $A'$, $B'$ 分别同构于 $A$, $B$. 提供一种切实制造这种同构副本的办法 (正如我们在 \S\ref{subsec:disjoint-unions-products} 中所做的) 是容易的; 但事实上这容易得有点过头---许多别的选择同样可能, 而且没有一个看起来比其他的更好. 因此, 更明智的做法其实是: 不要一劳永逸地作一个固定的选择, 而是干脆接受这样一个事实: 所有这些选择都给出 $A \amalg B$ 的合格候选者. 从这种平等主义的立场看, 运算 $A \amalg B$ 的结果作为上述意义下的集合并不``良定义''. 然而, 容易看出 (习题~\ref{ex:disjoint-union-welldef}) $A \amalg B$ 在同构意义下是良定义的: 也就是说, 副本 $A'$, $B'$ 的任何两种选取都导致同构的 $A \amalg B$ 候选者. 同样的考虑适用于积与商.

\index{fanchou@范畴(category)!Ab@Ab!G-Mod@G-Mod}由取无交并、积或商得到的集合, 其主要特征其实并不在于``它们包含哪些元素'', 而在于``它们与所有其他集合的关系''. 当我们在范畴的语境中重温这些运算以及其他运算%
\footnote{读者还应当知道, 我们迄今见过的这些运算还存在若干重要的变体---特别重要的是积与无交并的``纤维化''版本.}%
时, 这一点将 (更加) 清楚.

\subsection*{习题}
\phantomsection
\addcontentsline{toc}{subsection}{习题}

\begin{prob}
\label{ex:bijections-count}
$\triangleright$ 一个有 $n$ 个元素的集合 $S$ 到其自身有多少个不同的双射? [\S~II.\ref{subsec:symmetric-groups}]
\end{prob}

\begin{prob}
\label{ex:right-inverse}
$\triangleright$ 证明命题~\ref{prop:left-right-inverse} 中的断言 (2). 你可以假定: 给定一个集合的一族两两不交的非空子集, 存在一种办法从该族的每个成员中选取一个元素%
\footnote{这条 (合理的) 陈述就是选择公理; 参见 \S V.\ref{sec:zorn-lemma}.}%
. [\S\ref{subsec:second-viewpoint}, V.\ref{ex:ac-right-inverse}]
\end{prob}

\begin{prob}
\label{ex:inverse-bijection}
证明双射的逆是双射, 并且两个双射的复合是双射.
\end{prob}

\begin{prob}
\label{ex:iso-equiv-rel}
$\triangleright$ 证明``同构''是一个 (在任何由集合构成的集合上的) 等价关系. [\S\ref{subsec:isomorphisms}]
\end{prob}

\begin{prob}
\label{ex:epimorphism}
$\triangleright$ 按照 \S\ref{subsec:mono-epi} 中单态射概念的风格, 提出满态射的概念, 并就满态射与满射证明一个与命题~\ref{prop:mono-injective} 类似的结果. (提示: 翻转所有箭头.)\erratumfootnote{根据官方勘误, 此处补充提示内容. } [\S\ref{subsec:mono-epi}, \S\ref{subsec:mono-epi-cats}]
\end{prob}

\begin{prob}
\label{ex:section-of-projection}
采用例~\ref{ex:projections} 的记号, 解释任意映射 $f : A \to B$ 如何确定 $\pi_A$ 的一个截面.
\end{prob}

\begin{prob}
\label{ex:graph-iso}
设 $f : A \to B$ 是任意映射. 证明 $f$ 的图像 $\Gamma_f$ 同构于 $A$.
\end{prob}

\begin{prob}
\label{ex:canonical-decomp-exp}
尽可能显式地描述映射 $\R \to \CC$, $r \mapsto e^{2\pi i r}$ 的典范分解 (参见 \S\ref{subsec:canonical-decomposition}) 中的所有项. (这道题与此前布置过的一道题相配. 是哪一道?)
\end{prob}

\begin{prob}
\label{ex:disjoint-union-welldef}
$\triangleright$ 证明: 若 $A' \cong A''$ 且 $B' \cong B''$, 并且 $A' \cap B' = \emptyset$, $A'' \cap B'' = \emptyset$, 则 $A' \cup B' \cong A'' \cup B''$. 由此断定, 运算 $A \amalg B$ (如 \S\ref{subsec:disjoint-unions-products} 中所描述的) 在同构意义下是良定义的 (参见 \S\ref{subsec:clarification}). [\S\ref{subsec:clarification}, \ref{ex:redo-2-9}]
\end{prob}

\begin{prob}
\label{ex:card-ba}
$\triangleright$ 证明: 若 $A$ 与 $B$ 是有限集, 则 $|B^A| = |B|^{|A|}$. [\S\ref{subsec:functions-def}, \ref{ex:two-to-power}, \S~II.\ref{subsec:hom-examples}]
\end{prob}

\begin{prob}
\label{ex:two-to-power}
$\triangleright$ 鉴于习题~\ref{ex:card-ba}, 用 $2^A$ 表示从任意集合 $A$ 到一个二元集合 (比如 $\{0,1\}$) 的映射之集并非没有道理. 证明 $2^A$ 与 $A$ 的幂集之间存在双射 (参见 \S\ref{subsec:inclusion}). [\S\ref{subsec:inclusion}, III.\ref{ex:powerset-iso}]
\end{prob}

\section{范畴}
\label{sec:categories}

\index{fanchou@范畴(category)!Ab@Ab!Grp@Grp}范畴%

的语言被亲切地称为抽象废话%
\index{chouxinangfeihua@抽象废话(abstract nonsense)}
\index{Steenrod@Steenrod!Norman@Norman}\index{fanshu@范数(norm)!qianru@嵌入(embedding)!yukuozhangyuansude@域扩张元素的(of an element of a field extension)}\index{fanshu@范数(norm)!yukuozhangyuansude@域扩张元素的(of an element of a field extension)}\index{fanshu@范数(norm)!zai Z[i] zhong@在 Z[i] 中(in Z[i])}\index{fanshu@范数(norm)!zai Z[A] zhong@在 Z[A] 中(in Z[A])}\index{fanshu@范数(norm)!zai Z[√2] zhong@在 Z[√2] 中(in Z[√2])}\index{fanshu@范数(norm)!zai Q(√d) zhong@在 Q(√d) 中(in Q(√d))}\index{fanshu@范数(norm)!zai R-Mod zhong@在 R-Mod 中(in R-Mod)}\index{fanshu@范数(norm)}, 这个名字是 Norman Steenrod 起的. 这一称呼基本上是准确的, 而且未必带贬义: 范畴之所以关乎``废话''\,(nonsense), 是因为它们只关心所表示对象的``结构'', 而不关心其``意义''. 重点不在于你是如何碰到手头这个特定的集合的, 而在于这个集合如何与所有其他集合发生关系. 更糟 (或更好) 的是, 重点不在于研究集合以及集合间的映射, 而在于研究``东西, 以及从一个东西通向另一个东西的东西'', 而不必明确说出这些东西是什么: 它们可以是集合, 也可以是群、环、向量空间、模, 或者是其他如此另类的东西, 以至于读者 (眼下) 根本无权了解它们.

\index{fanchou@范畴(category)!Ab@Ab!G-Set@G-Set}``范畴''初看之下会在多方面让人联想到集合. 范畴可能让你想起集合, 因为它们是``对象的汇集'', 并且还会有``从范畴到范畴的映射'' (称为函子%
\index{hanzide@函子的(of functors)}\footnote{不过, 我们要到后面的章节才会考虑函子: 我们与函子的第一次正式会面将在第八章.}%
\index{hanzi@函子(functor)}
\index{fanchou@范畴(category)!Ab@Ab!k-Vect@k-Vect}). 同时, 每个范畴又可能让你想起全体集合的汇集, 因为在它所包含的东西之间会有``映射''的类似物.

\subsection{定义}
\label{subsec:categories-def}

\index{fanchou@范畴(category)!Ab@Ab!R-Alg@R-Alg}范畴的定义初看很复杂, 但其要义可以很快概括如下: 一个范畴由一个``对象''的汇集以及这些对象之间的``态射''%
\index{taishe@态射(morphism)}
\index{ziran@自然(natural)}构成, 并满足一列自然的条件.

\index{fanchou@范畴(category)!Ab@Ab!huan@环(ring)}读者会注意到, 我避免写``一个对象的集合'', 而选用更一般的``汇集''. 这是一个令人恼火却又无法回避的困难: 例如, 我们想要一个``集合的范畴'', 其中``对象''是集合, ``态射''是集合间的映射, 问题在于, 根本不存在一个由所有集合组成的集合%
\footnote{这是我们从罗素悖论中学到的事情之一.}%
. 从某种意义上说, 全体集合的汇集``太大'', 而无法成为一个集合. 不过, 处理这类``汇集''的办法是有的, 它们的技术名称是类%
\index{lei@类(class)}
\index{qun@群(group)!zuoyong@作用(action)}. 存在一个由所有集合组成的``类'' (并且也将有照管群、环等等的类).

\index{zugou@足够(enough)}另一种替代方案是定义一个足够大的集合 (称为宇宙%
\index{yuzhou@宇宙(universe)}
\index{fanchou@范畴(category)!Ab@Ab!R-Mod@R-Mod}), 然后约定所有范畴的所有对象都将从这个庞然大物中选取.

\index{fanchou@范畴(category)!Ab@Ab!jihe@集合(set)}无论如何, 读者对此需要知道的全部就是: 有办法让它行得通. 我们将在定义中使用``类''这个词, 但这不会影响本书中的任何证明或任何其他定义. 此外, 在下面考虑的某些例子中, 相关的类是一个集合 (此时我们称该范畴是小%
\index{xiaofanchou@小范畴(small category)}
的), 所以读者在面对这些例子时会有宾至如归之感.

\begin{definition}
\index{fanchou@范畴(category)!Ab@Ab!daiji@带积(with products)}\index{fanchou@范畴(category)!Ab@Ab!daiyuji@带余积(with coproducts)}\index{fanchou@范畴(category)!Ab@Ab!sanjiao@三角(triangle)}\index{fanchou@范畴(category)!Ab@Ab!zifanchou@子范畴(subcategory)}\index{fanchou@范畴(category)!Ab@Ab!xiao@小(small)}\index{fanchou@范畴(category)!Ab@Ab!qiepian@切片(slice)}\index{fanchou@范畴(category)!Ab@Ab!yujiaxing@预加性(preadditive)}\index{fanchou@范畴(category)!Ab@Ab!fan@反(opposite)}\index{fanchou@范畴(category)!Ab@Ab!hanzide@函子的(of functors)}\index{fanchou@范畴(category)!Ab@Ab!fuxingde@复形的(of a complex)}\index{fanchou@范畴(category)!Ab@Ab!fangshedaishujide@仿射代数集的(of affine algebraic sets)}\index{fanchou@范畴(category)!Ab@Ab!dedengjia@的等价(equivalence of)}\index{fanchou@范畴(category)!Ab@Ab!lisan@离散(discrete)}\index{fanchou@范畴(category)!Ab@Ab!daochu@导出(derived)}\index{fanchou@范畴(category)!Ab@Ab!yuqiepian@余切片(coslice)}\index{fanchou@范畴(category)!Ab@Ab!douhao@逗号(comma)}\index{fanchou@范畴(category)!Ab@Ab!jiafa@加法(addition)}\index{fanchou@范畴(category)!Ab@Ab!Abel@Abel}\index{fanchou@范畴(category)!Ab@Ab!Seq(A)@Seq(A)}\index{fanchou@范畴(category)!Ab@Ab!Set*@Set*}\label{def:category}
一个范畴 $\cate{C}$ 由如下数据构成:
\begin{itemize}
	\item 一个由范畴的对象构成的类 $\Obj(\cate{C})$;
	\item 以及, 对 $\cate{C}$ 的每两个对象 $A$, $B$, 一个由态射构成的集合 $\Hom_{\cate{C}}(A, B)$, 它们具有下面列出的性质.\index[sym1]{25@$\Hom_{\cate{C}}(A,B)$}
\end{itemize}
\end{definition}

\index{ziran@自然(natural)!tonggou@同构(isomorphism)}要记住的原型是: 把对象想成``集合'', 把态射想成``映射''. 仅此一个例子就应当使态射的定义性质显得自然且容易记忆:
\begin{itemize}
	\item 对 $\cate{C}$ 的每个对象 $A$, 存在 (至少) 一个态射 $1_A \in \Hom_{\cate{C}}(A, A)$, 即 $A$ 上的``恒同''.
	\item 态射可以复合: 两个态射 $f \in \Hom_{\cate{C}}(A, B)$ 与 $g \in \Hom_{\cate{C}}(B, C)$ 确定一个态射 $gf \in \Hom_{\cate{C}}(A, C)$. 也就是说, 对 $\cate{C}$ 中对象的每个三元组 $A$, $B$, $C$, 有一个 (集合的) 映射
	\[
	\Hom_{\cate{C}}(A, B) \times \Hom_{\cate{C}}(B, C) \to \Hom_{\cate{C}}(A, C),
	\]
	并且 $(f,g)$ 的像记作 $gf$.
	\item 这个``复合律''是结合的: 若 $f \in \Hom_{\cate{C}}(A, B)$, $g \in \Hom_{\cate{C}}(B, C)$, $h \in \Hom_{\cate{C}}(C, D)$, 则
	\[
	(hg)f = h(gf).
	\]
	\item 恒同态射关于复合是恒同的: 也就是说, 对所有 $f \in \Hom_{\cate{C}}(A, B)$, 有
	\[
	f 1_A = f, \qquad 1_B f = f.
	\]
\end{itemize}
这确实有点儿拗口, 但再说一遍, 要记住这一切, 只需想想集合的映射. 还有一条进一步的要求: 集合
\[
\Hom_{\cate{C}}(A, B), \qquad \Hom_{\cate{C}}(C, D)
\]
在 $A \neq C$ 或 $B \neq D$ 时不相交; 这是一件你通常不会去想的事情, 但它对普通的集合映射同样成立%
\footnote{我们会常用``集合映射''(set-function)一词, 以强调我们是在集合的语境中谈论一个映射.}%
. 也就是说, 若两个映射是同一个映射, 则它们必然有相同的源和相同的目标: 源与目标是集合映射的数据的一部分.

范畴 $\cate{C}$ 的对象 $A$ 到其自身的态射称为自同态%
\index{zitongtai@自同态(endomorphism)}
; $\Hom_{\cate{C}}(A, A)$ 记作 $\End_{\cate{C}}(A)$. 范畴的一条公理告诉我们这是一个``带基点的''集合, 因为 $1_A \in \End_{\cate{C}}(A)$. 读者应当注意, 复合在 $\End_{\cate{C}}(A)$ 上定义了一个``运算'': 若 $f$, $g$ 是 $\End_{\cate{C}}(A)$ 的元素, 则它们的复合 $gf$ 亦然.\index[sym1]{14@$\End_{\cate{C}}(A)$}

长期写下``$f \in \Hom_{\cate{C}}(A, B)$''会让人厌倦. 如果范畴已在上下文中明确, 便可以放心地省去下标 $\cate{C}$, 甚至像对集合映射那样使用箭头: $f : A \to B$. 这也使得我们可以在任何范畴中画态射的图; 若沿图的所有路径复合态射都得到相同的结果---正如 \S\ref{subsec:composition} 中对集合映射之图的解释那样---则称该图``交换'' (或是``交换图'').

事实上, 我们今后将放心地把图当作范畴中可能的候选对象. 在此语境下, 图的正式定义是: 范畴 $\cate{C}$ 的一组对象, 连同这些对象之间规定的一些态射; 若图在上述意义下交换, 则称图交换. 图的视觉呈现的具体样式自然是无关紧要的.

\subsection{例子}
\label{subsec:categories-ex}

读者应当注意, 范畴概念的定义有 $90\%$ 花在解释其态射的性质上; 可以公平地说, 态射才是范畴中重要的组成成分. 尽管如此, 从心理上说, 人们还是难以抗拒地要用对象来思考范畴: 例如, 人们会说``集合的范畴''. 关键在于, 通常所能考虑的``态射''的种类是 (至少在心理上) 由对象决定的: 如果谈的是集合, 除了集合的映射之外, ``态射''还能是什么意思呢? 在别的情况 (参见下面的例~\ref{ex:slice-categories} 或习题~\ref{ex:mset-category}) 中, 态射应当是什么就不那么清楚了, 而寻找``正确的''概念可能是一个有趣的课题.

\begin{example}
\label{ex:set-category}
到现在为止, ``集合 (作为对象) 连同集合映射 (作为态射) 构成一个范畴''应当是再清楚不过的了; 若否, 读者必须止步于此, 直到这一断言不再残留任何神秘之处为止%
\footnote{我们会时不时给读者这样的提示: 在关键时刻, 盘点一下自己已经知道了什么, 比蒙头向前赶路更有用. 此刻感到困难, 说明需要仔细重读前面的内容. 若谜团依旧, 办公答疑时间就是为此而设的. 不过, 通常你应当能够凭借我们给出的信息自己找到出路, 而且这样做你多半会学到更多. 在寻求专业帮助之前, 请先尽自己最大的努力.}%
.

这个重要的范畴没有获得普遍接受的官方记号. 习惯上写作``Set''或``Sets'', 并配以某种花哨的装饰以示强调. 例如, 在文献中人们可以遇到 SET, Sets, Set, (Sets), 以及这些主题的许多有趣的变体. 我们将用``无衬线''字体表示范畴; 于是, $\cate{Set}$ 表示集合的范畴. 这样
\begin{itemize}
	\item $\Obj(\cate{Set}) =$ 全体集合构成的类;
	\item 对 $\Obj(\cate{Set})$ 中的 $A$, $B$ (即 $A$, $B$ 是集合), $\Hom_{\cate{Set}}(A, B) = B^A$.
\end{itemize}
注意, \S\ref{subsec:operations}--\S\ref{subsec:equivalence-relations} 中回顾的那些运算的存在并不是范畴定义的一部分: 这些运算凸显了 $\cate{Set}$ 的一些有趣特征, 其他范畴可以有也可以没有. 我们很快会回到其中某些运算, 更精确地理解它们关于 $\cate{Set}$ 说明了什么.
\end{example}

\begin{example}
\label{ex:relation-category}
这里是一个截然不同的例子.

设 $S$ 是一个集合, $\sim$ 是 $S$ 上一个满足自反性与传递性的关系. 那么我们可以把这些数据编码成一个范畴:
\begin{itemize}
	\item 对象: $S$ 的元素;
	\item 态射: 若 $a$, $b$ 是对象 (即 $a, b \in S$), 则当 $a \sim b$ 时令 $\Hom(a,b)$ 为由元素 $(a,b) \in S \times S$ 构成的集合, 否则令 $\Hom(a,b) = \emptyset$.
\end{itemize}
注意 (与 $\cate{Set}$ 不同) 这里的态射非常少: 每一对对象至多一个, 而``没有关系''的对象之间根本没有态射.

我们必须定义``态射的复合'', 并验证 \S\ref{subsec:categories-def} 中规定的条件得到满足. 首先, 有``恒同''吗? 若 $a$ 是一个对象 (即 $a \in S$), 我们需要找到一个元素
\[
1_a \in \Hom(a, a).
\]
这恰恰是我们要假定 $\sim$ 自反的原因: 它告诉我们 $\forall a,\ a \sim a$; 也就是说, $\Hom(a,a)$ 由单个元素 $(a,a)$ 构成. 于是我们别无选择: 只能让
\[
1_a = (a, a) \in \Hom(a, a).
\]
至于复合, 设 $a, b, c$ 是对象 (即 $S$ 的元素), 且
\[
f \in \Hom(a, b), \quad g \in \Hom(b, c);
\]
我们必须定义相应的态射 $gf \in \Hom(a, c)$. 现在,
\[
f \in \Hom(a, b)
\]
告诉我们 $\Hom(a,b)$ 非空, 而按照这个范畴中态射的定义, 这意味着 $a \sim b$, 且 $f$ 正是 $S \times S$ 中的元素 $(a,b)$. 类似地, $g \in \Hom(b,c)$ 告诉我们 $b \sim c$ 且 $g = (b,c)$. 由于我们假定 $\sim$ 传递,
\[
a \sim b \text{ 并且 } b \sim c \implies a \sim c.
\]
这告诉我们 $\Hom(a,c)$ 由单个元素 $(a,c)$ 构成. 于是我们又别无选择: 只能让
\[
gf := (a, c) \in \Hom(a, c).
\]
这个运算是结合的吗? 若 $f \in \Hom(a, b)$, $g \in \Hom(b, c)$, $h \in \Hom(c, d)$, 则必然有
\[
f = (a, b), \quad g = (b, c), \quad h = (c, d)
\]
以及
\[
gf = (a, c), \quad hg = (b, d),
\]
从而
\[
h(gf) = (a, d) = (hg)f,
\]
结合性得证.

读者将不难验证 $1_a$ 关于这个复合是恒同, 正如所需 (习题~\ref{ex:identity-relation-category}).

这个构造最平凡的实例, 是取集合 $S$ 配上等价关系``=''所得到的范畴; 也就是说, 唯一的态射是恒同态射. 这样的范畴称为离散范畴%
\index{lisancate@离散范畴(discrete category)}
.

作为另一个例子, 考虑给 $\Z$ 配上关系 $\leq$ 所对应的范畴. 例如,
\begin{center}
\begin{tikzcd}[row sep=2.6em, column sep=1.9em]
	2 \ar[r] \ar[rd] & 3 \ar[r] \ar[d, "1_3"'] & 5 \ar[r] & 7 \\
	& 3 \ar[r] \ar[urr] & 4 \ar[ru] &
\end{tikzcd}
\end{center}
是这个范畴中一个 (随意选取的) 交换图. 即便把竖直箭头 $3 \to 3$ 反向, 或者添加一条从 $3$ 到 $4$ 的箭头, 它仍将是这个范畴中的 (交换) 图; 但我们不允许画一条从 $4$ 到 $3$ 的箭头, 因为 $4 \not\leq 3$.

这些范畴非常特殊---例如, 在其中画的每个图都必然交换, 而在比如 $\cate{Set}$ 中情况远非如此. 还要注意, 这些范畴都是小范畴.
\end{example}

\begin{example}
\label{ex:power-set-category}
这里又有另一个%
\footnote{事实上, 这又是例~\ref{ex:relation-category} 中所考虑的那类范畴的一个实例. 你看出为什么了吗? (习题~\ref{ex:power-set-as-relation}.)}%
小范畴的例子.

仍设 $S$ 是一个集合. 如下定义范畴 $\widehat{S}$:
\begin{itemize}
	\item $\Obj(\widehat{S}) = \mathscr{P}(S)$, 即 $S$ 的幂集 (参见 \S\ref{subsec:inclusion} 与习题~\ref{ex:two-to-power});
	\item 对 $\widehat{S}$ 的对象 $A$, $B$ (即 $A \subseteq S$ 且 $B \subseteq S$), 当 $A \subseteq B$ 时令 $\Hom_{\widehat{S}}(A, B)$ 为由序对 $(A,B)$ 构成的集合\erratumfootnote{根据官方勘误, 此处原书作``为序对 $(A,B)$'', 态射集应为由该序对构成的单元素集.}, 否则令 $\Hom_{\widehat{S}}(A, B) = \emptyset$.
\end{itemize}
恒同 $1_A$ 由序对 $(A,A)$ 构成 (它是 $A$ 到 $A$ 的一个、事实上也是唯一的一个态射, 因为 $A \subseteq A$). 复合由串接包含关系得到: 若 $\widehat{S}$ 中有态射
\[
A \to B, \quad B \to C,
\]
则 $A \subseteq B$ 且 $B \subseteq C$; 从而 $A \subseteq C$, 于是有态射 $A \to C$. 检验 \S\ref{subsec:categories-def} 规定的各条公理应当是例行公事 (务必确认确属如此!).

这种风格的例子 (但采用更精致的结构, 比如拓扑空间的开子集族) 在代数几何这样的成熟领域中极为重要.
\end{example}

\begin{example}
\label{ex:slice-categories}
下一个例子非常抽象, 但对它的思考将使你对我们迄今所见的一切感到相当自在; 而且它是一种非常常见的构造, 其各种变体将在本课程中大量涌现.

设 $\cate{C}$ 是一个范畴, $A$ 是 $\cate{C}$ 的一个对象. 我们要定义一个范畴 $\cate{C}_A$, 其对象是 $\cate{C}$ 中的某些态射, 而其态射是 $\cate{C}$ 中的某些图 (想不到吧!).

\begin{itemize}
	\item $\Obj(\cate{C}_A) =$ 从 $\cate{C}$ 的任意对象到 $A$ 的所有态射; 于是, $\cate{C}_A$ 的对象是对 $\cate{C}$ 的某个对象 $Z$ 而言的态射 $f \in \Hom_{\cate{C}}(Z, A)$. 用图形说, $\cate{C}_A$ 的对象是 $\cate{C}$ 中一个箭头 $Z \xrightarrow{f} A$; 这些常被画成``自上而下''的样子, 如
	\begin{center}
	\begin{tikzcd}[row sep=1.7em]
		Z \ar[d, "f"'] \\
		A
	\end{tikzcd}
	\end{center}

	$\cate{C}_A$ 的态射应当是什么呢? 对一个以上述对象为对象的范畴指定态射, 实际上只有一种合理的方式. 勇敢的读者会愿意就此停止阅读, 在独立想出定义之后再继续. 在本书将遇到的各种构造背后潜伏着许多类似的例子; 理想情况下, 当时机到来时它们应当显得完全自然. 现在花一点力气理解这个原型情形, 将来必有丰厚回报. 剧透在后面, 所以请现在就把这份讲义放到一边, 勾勒出 $\cate{C}_A$ 中态射的定义.

	欢迎回来.

	\item 设 $f_1$, $f_2$ 是 $\cate{C}_A$ 的对象, 即 $\cate{C}$ 中的两个箭头
	\begin{center}
	\begin{tikzcd}[row sep=1.7em, column sep=2.5em]
		Z_1 \ar[d, "f_1"'] & & Z_2 \ar[d, "f_2"'] \\
		A & & A
	\end{tikzcd}
	\end{center}
	态射 $f_1 \to f_2$ 定义为``环境''范畴 $\cate{C}$ 中的交换图
	\begin{center}
	\begin{tikzcd}
		& A & \\
		Z_1 \ar[ur, "f_1"] \ar[rr, "\sigma"'] & & Z_2 \ar[ul, "f_2"']
	\end{tikzcd}
	\end{center}
\end{itemize}
也就是说, 态射 $f_1 \to f_2$ 恰好对应于 $\cate{C}$ 中那些使 $f_1 = f_2 \sigma$ 的态射 $\sigma : Z_1 \to Z_2$.

一旦明白了态射必须是什么, 检验它们满足 \S\ref{subsec:categories-def} 中列出的公理便是直截了当的. 恒同从 $\cate{C}$ 的恒同继承而来: 对 $\cate{C}_A$ 中的 $f : Z \to A$, 恒同 $1_f$ 对应于图
\begin{center}
\begin{tikzcd}
	Z \ar[rr, "1_Z"] \ar[rd, "f"'] & & Z \ar[ld, "f"] \\
	& A &
\end{tikzcd}
\end{center}
它交换, 因为 $\cate{C}$ 是范畴. 复合同样是 $\cate{C}$ 中复合的副产品. $\cate{C}_A$ 中两个态射 $f_1 \to f_2 \to f_3$ 对应于把两个交换图并排放置:
\begin{center}
\begin{tikzcd}
	Z_1 \ar[r, "\sigma"] \ar[rd, "f_1"'] & Z_2 \ar[r, "\tau"] \ar[d, "f_2"'] & Z_3 \ar[dl, "f_3"] \\
	& A &
\end{tikzcd}
\end{center}
于是 (又因为 $\cate{C}$ 是范畴!) 由此得到: 去掉中间箭头所得的图, 即
\begin{center}
\begin{tikzcd}
	Z_1 \ar[rr, "\tau\sigma"] \ar[rd, "f_1"'] & & Z_3 \ar[ld, "f_3"] \\
	& A &
\end{tikzcd}
\end{center}
也交换. 把这一切都检验一遍 (!), 并验证 $\cate{C}_A$ 中的复合是结合的 (同样, 这由 $\cate{C}$ 中复合的结合性立即得到).

按这种方式构造的范畴在文献中称为切片范畴%
\index{qiepianfanchou@切片范畴(slice category)}
; 它们是逗号范畴%
\index{douhaofanchou@逗号范畴(comma category)}
的特别情形.
\end{example}

\begin{example}
\label{ex:slice-concrete}
为了具体起见, 让我们把例~\ref{ex:slice-categories} 中的构造应用到例~\ref{ex:relation-category} 所构造的范畴上, 比如取 $S = \Z$, $\sim$ 为关系 $\leq$. 称这个范畴为 $\cate{C}$, 并选取 $\cate{C}$ 的一个对象 $A$---即一个整数, 例如 $A = 3$. 那么 $\cate{C}_A$ 的对象是 $\cate{C}$ 中以 $3$ 为目标的态射, 即满足 $n \leq 3$ 的序对 $(n,3) \in \Z \times \Z$. 存在态射
\[
(m, 3) \to (n, 3)
\]
当且仅当 $m \leq n$. 在这个例子中, 可以无伤大雅地把 $\cate{C}_A$ 认同于整数 $\leq 3$ 的``子范畴'', 其态射与 $\cate{C}$ 中``相同''.
\end{example}

\begin{example}
\label{ex:coslice}
考虑 $\cate{C}$ 中从固定对象 $A$ 出发到 $\cate{C}$ 的所有对象的态射, 同样用适当的交换图定义态射, 便得到一个与例~\ref{ex:slice-categories} 所探索的例子完全类似的例子. 这引出余切片范畴%
\index{yuqiepianfanchou@余切片范畴(coslice category)}
. 该构造的细节由读者补出 (习题~\ref{ex:coslice-details}).
\end{example}

\begin{example}
\label{ex:pointed-sets}
作为例~\ref{ex:coslice} 中那种范畴的一个``具体''实例, 取 $\cate{C} = \cate{Set}$, $A =$ 一个固定的单元素集 $\{*\}$. 把所得的范畴称为 $\cate{Set}^{*}%
$.

于是, $\cate{Set}^{*}$ 中的一个对象是 $\cate{Set}$ 中的一个态射 $f : \{*\} \to S$, 其中 $S$ 是任意集合. 因此, $\cate{Set}^{*}$ 中一个对象的信息由一个非空集合 $S$ 的选取以及一个元素 $s \in S$ 的选取构成---即元素 $f(*)$: 这个元素决定 $f$, 也由 $f$ 决定.

于是, 我们可以把 $\cate{Set}^{*}$ 的对象记作 $(S, s)$, 其中 $S$ 是任意集合, $s \in S$ 是 $S$ 的任意元素.

两个这样的对象之间的态射 $(S, s) \to (T, t)$, 则对应于 (请自行验证!) 一个满足 $\sigma(s) = t$ 的集合映射 $\sigma : S \to T$.

$\cate{Set}^{*}$ 的对象称为带基点的集合%
\index{daijidiandejih@带基点的集合(pointed set)}
. 我们在本书中将研究的许多结构都会是带基点的集合. 例如 (我们将会看到), 一个``群''是一个集合 $G$, 在其他要求之外还带有一个特定元素 $e_G$ (它的``恒同元''); ``群同态''将是那些在别的性质之外还把恒同元送到恒同元的映射; 于是, 它们就是上述意义下带基点的集合的态射.
\end{example}

\begin{example}
\label{ex:two-objects}
按照例~\ref{ex:slice-categories} 与例~\ref{ex:coslice} 的风格再思考几个``抽象''的例子是有益的. 它们将是前文许诺过的、对 \S\ref{subsec:operations} 中提到的某些运算的重新审视中的本质要素. 对于已经掌握了例~\ref{ex:slice-categories} 与例~\ref{ex:coslice} 的读者来说, 这些定义会显得朴素得令人失望.

这一次我们从给定的范畴 $\cate{C}$ 及其两个对象 $A$, $B$ 出发. 我们可以按照定义 $\cate{C}_A$ 时所用的本质相同的程序, 定义一个新范畴 $\cate{C}_{A,B}$:
\begin{itemize}
	\item $\Obj(\cate{C}_{A,B}) = \cate{C}$ 中的图
	\begin{center}
	\begin{tikzcd}
		& A & \\
		Z \ar[ur, "f"] \ar[dr, "g"'] & & \\
		& B &
	\end{tikzcd}
	\end{center}
	\item 以及态射
	\begin{center}
	\begin{tikzcd}
		& A & & & & A & \\
		Z_1 \ar[ur, "f_1"] \ar[dr, "g_1"'] & & \longrightarrow & & Z_2 \ar[ur, "f_2"] \ar[dr, "g_2"'] & & \\
		& B & & & & B &
	\end{tikzcd}
	\end{center}
	即交换图
	\begin{center}
	\begin{tikzcd}[column sep=2.2em, row sep=1.2em]
		& & A & \\
		Z_1 \ar[urr, bend left=16, "f_1"] \ar[r, "\sigma"'] \ar[drr, bend right=16, "g_1"'] & Z_2 \ar[ur, "f_2"] \ar[dr, "g_2"'] & & \\
		& & B &
	\end{tikzcd}
	\end{center}
\end{itemize}
把这个粗略的描述形式化的任务留给读者. 这个例子确实不过是 $\cate{C}_A$ 与 $\cate{C}_B$ 的混合体, 两种结构之所以相互作用, 是由于同一个 $\sigma$ 必须使图的两侧同时交换这一严格要求:
\[
f_1 = f_2 \sigma \quad \text{并且} \quad g_1 = g_2 \sigma,
\]
``同时''成立.

把大部分箭头翻转, 就得到例~\ref{ex:coslice} 的一个类似的变体, 产生一个范畴, 我们可以%
\footnote{这些常见的范畴似乎并没有公认的记号.}%
把它记作 $\cate{C}^{A,B}$; 细节留给读者.
\end{example}

\begin{example}
\label{ex:fibered}
作为这些例子的最后一个变体, 我们以考虑 $\cate{C}_{A,B}$ (以及 $\cate{C}^{A,B}$) 的纤维化版本收尾. 把它当作一个检验, 看看你是否真正理解了 $\cate{C}_{A,B}$---专家会告诉你, 对刚刚学习范畴的学生来说这看起来相当高深, 所以如果起初不够顺畅也不要灰心 (但如果很顺畅, 请拍拍自己的肩膀!). 从给定的范畴 $\cate{C}$ 出发, 这次在 $\cate{C}$ 中选取两个有公共目标 $C$ 的固定态射 $\alpha : A \to C$, $\beta : B \to C$. 于是我们可以如下考虑一个范畴 $\cate{C}_{\alpha,\beta}$:
\begin{itemize}
	\item $\Obj(\cate{C}_{\alpha,\beta}) = \cate{C}$ 中的交换图
	\begin{center}
	\begin{tikzcd}[row sep=1.2em, column sep=2.2em]
		& A \ar[dr, "\alpha"] & \\
		Z \ar[ur, "f"] \ar[dr, "g"'] & & C \\
		& B \ar[ur, "\beta"'] &
	\end{tikzcd}
	\end{center}
	\item 以及对应于交换图的态射
	\begin{center}
	\begin{tikzcd}[column sep=1.9em, row sep=1.1em]
		& & A \ar[dr, "\alpha"] & \\
		Z_1 \ar[urr, bend left=14, "f_1"] \ar[r, "\sigma"'] \ar[drr, bend right=14, "g_1"'] & Z_2 \ar[ur, "f_2"] \ar[dr, "g_2"'] & & C \\
		& & B \ar[ur, "\beta"'] &
	\end{tikzcd}
	\end{center}
\end{itemize}
对例~\ref{ex:two-objects} 的扎实理解会使这个例子看起来同样驯顺; 此时读者把形式化的工作 (即解释复合如何进行、恒同是什么等等) 完成应当毫无困难.

同样留给读者的还有``镜像''例子 $\cate{C}^{\alpha,\beta}$ 的构造, 它从两个有公共源 $C$ 的态射 $\alpha : C \to A$, $\beta : C \to B$ 出发.
\end{example}

\subsection*{习题}
\phantomsection
\addcontentsline{toc}{subsection}{习题}

\begin{prob}
\label{ex:opposite-category}
$\triangleright$ 设 $\cate{C}$ 是一个范畴. 考虑一个结构 $\cate{C}\op$, 规定
\begin{itemize}
	\item $\Obj(\cate{C}\op) := \Obj(\cate{C})$;
	\item 对 $\cate{C}\op$ 的对象 $A$, $B$ (因而也是 $\cate{C}$ 的对象), $\Hom_{\cate{C}\op}(A, B) := \Hom_{\cate{C}}(B, A)$.
\end{itemize}
说明如何把它做成一个范畴 (即定义 $\cate{C}\op$ 中态射的复合, 并验证 \S\ref{subsec:categories-def} 中列出的各条性质).

直观上, ``反''范畴%
\index{fanfanchou@反范畴(opposite category)}
$\cate{C}\op$ 不过是把 $\cate{C}$ 中``所有箭头反转''所得. [\ref{ex:final-initial-op}, \S~VIII.\ref{subsec:functors}, \S~IX.\ref{subsec:additive-categories}, IX.\ref{ex:1.10}]
\end{prob}

\begin{prob}
\label{ex:end-finite}
若 $A$ 是有限集, $\End_{\cate{Set}}(A)$ 有多大?
\end{prob}

\begin{prob}
\label{ex:identity-relation-category}
$\triangleright$ 精确表述``$1_a$ 是例~\ref{ex:relation-category} 中复合的恒同''的含义, 并证明这一断言. [\S\ref{subsec:categories-ex}]
\end{prob}

\begin{prob}
\label{ex:lt-category}
我们能按例~\ref{ex:relation-category} 的风格用 $\Z$ 上的关系 $<$ 定义一个范畴吗?
\end{prob}

\begin{prob}
\label{ex:power-set-as-relation}
$\triangleright$ 解释在何种意义下例~\ref{ex:power-set-category} 是例~\ref{ex:relation-category} 中所考虑的那些范畴的一个实例. [\S\ref{subsec:categories-ex}]
\end{prob}

\begin{prob}
\label{ex:matrix-category}
$\triangleright$ (假定对线性代数有一定的熟悉.) 按如下方式定义范畴 $\cate{V}$: 取 $\Obj(\cate{V}) = \N$, 并对所有 $n, m \in \N$ 令 $\Hom_{\cate{V}}(n, m) =$ 实元素的 $m \times n$ 矩阵之集. (行数或列数为 $0$ 的矩阵是什么意思, 留给读者自行领会.) 用矩阵的乘积定义复合. 这个范畴是否让你感到``似曾相识''? [\S~VI.\ref{subsec:matrices}, \S~VIII.\ref{subsec:equiv-categories}]
\end{prob}

\begin{prob}
\label{ex:coslice-details}
$\triangleright$ 仔细定义例~\ref{ex:coslice} 中的对象与态射, 并画出对应于复合的图. [\S\ref{subsec:categories-ex}]
\end{prob}

\begin{prob}
\label{ex:subcategory}
$\triangleright$ 范畴 $\cate{C}$ 的子范畴%
\index{zifanchou@子范畴(subcategory)}
$\cate{C}'$ 由 $\cate{C}$ 的一些对象组成, 并对 $\Obj(\cate{C}')$ 中的一切对象 $A$, $B$ 配备态射集 $\Hom_{\cate{C}'}(A,B) \subseteq \Hom_{\cate{C}}(A,B)$, 使得 $\cate{C}$ 中的恒同与复合使 $\cate{C}'$ 成为一个范畴. 子范畴 $\cate{C}'$ 称为全的%
\index{quande@全的(full)}
, 若对所有 $A, B \in \Obj(\cate{C}')$ 有 $\Hom_{\cate{C}'}(A,B) = \Hom_{\cate{C}}(A,B)$. 构造一个由无限集构成的范畴, 并解释如何把它看成 $\cate{Set}$ 的一个全子范畴. [\ref{ex:mono-composition}, \S~VI.\ref{subsec:r-mod}, \S~VIII.\ref{subsec:equiv-categories}]
\end{prob}

\begin{prob}
\label{ex:mset-category}
$\triangleright$ 对 \S\ref{subsec:multisets-indexed} 中引入的多重集概念的另一种替代, 是考虑配备了等价关系的集合, 其中等价的元素被取作``同一类''元素的多个实例. 定义这类增强了的结构之间的态射概念, 得到一个包含 $\cate{Set}$ (的一个``副本'') 作为全子范畴的范畴 $\cate{MSet}$. (可能有不止一种合理的做法! 这是一道有意留为开放式的习题.) $\cate{MSet}$ 中哪些对象按 \S\ref{subsec:multisets-indexed} 的定义确定普通的多重集? 如何确定? 从这个观点出发, 阐明多重集之间的态射应当是什么. (多重集的态射有几种自然的概念. 试着定义 $\cate{MSet}$ 中的态射, 使得你对普通多重集所得到的态射概念抓住你对这些对象的直观理解.) [\S\ref{subsec:multisets-indexed}, \S\ref{subsec:categories-ex}, \ref{ex:mono-epi-mset}]
\end{prob}

\begin{prob}
\label{ex:subobject-classifier}
\index{ziduixiangfenleizi@子对象分类子(subobject classifier)}由于范畴 $\cate{C}$ 的对象 (未必) 是集合, 一般说来如何通过推广子集的概念\erratumfootnote{根据官方勘误, 此处补充``通过推广子集的概念''(extrapolating the notion of subset).}以理解``子对象''并不清楚. 在某些情形下谈论子对象确实有意义, 并且 $\cate{C}$ 中任一给定对象 $A$ 的子对象与态射 $A \to \Omega$ 一一对应, 其中 $\Omega$ 是 $\cate{C}$ 的一个固定的特殊对象, 称为子对象分类器%
\index{ziduixiangfenleiqi@子对象分类器(subobject classifier)}
. 证明 $\cate{Set}$ 有子对象分类器.
\end{prob}

\begin{prob}
\label{ex:dual-two-objects}
$\triangleright$ 画出相关的图, 并为例~\ref{ex:two-objects} 中提到的范畴 $\cate{C}^{A,B}$ 定义复合与恒同. 对例~\ref{ex:fibered} 中提到的范畴 $\cate{C}^{\alpha,\beta}$ 做同样的事情. [\S\ref{subsec:coproducts-up}, \ref{ex:fibered-products}]
\end{prob}

\section{态射}
\label{sec:morphisms}

正如在 $\cate{Set}$ 中我们突出标出几类特殊的映射 (单射、满射、双射) 那样, 对任意范畴中的态射做同样的事情也是有用的. 读者应当注意, 在一般框架下, 通过态射对``元素''的作用来定义态射的品质是行不通的, 因为任意范畴的对象 (一般) 没有``元素''.

这正是我们此前花费一些时间从不同视角分析单性等概念的原因 (\S\ref{subsec:inj-surj-bij}--\S\ref{subsec:mono-epi}). 事实证明, 这些概念的另外几种视角确实可以顺利地移植到范畴的框架中.

\subsection{同构}
\label{subsec:isomorphisms}

设 $\cate{C}$ 是一个范畴.

\begin{definition}
\label{def:isomorphism}
态射 $f \in \Hom_{\cate{C}}(A, B)$ 称为同构%

, 如果它在复合下有 (双侧的) 逆: 也就是说, 若存在 $g \in \Hom_{\cate{C}}(B, A)$ 使得
\[
gf = 1_A, \qquad fg = 1_B.
\]
\end{definition}

回顾 \S\ref{subsec:second-viewpoint} 中集合双射 $f$ 的逆是``按元素''定义的; 特别地, 其定义没有任何歧义, 而且我们为该函数引入了记号 $f^{-1}$. 相比之下, 定义~\ref{def:isomorphism} 给出的``逆'' $g$, 其定义似乎没有显式地内建唯一性. 幸运的是, 它的定义性质确实保证了唯一性, 但这需要一番验证:

\begin{proposition}
\label{prop:inverse-unique}
同构的逆是唯一的.
\end{proposition}

\begin{proof}
我们必须验证: 若 $g_1$ 与 $g_2 : B \to A$ 都充当给定同构 $f : A \to B$ 的逆, 则 $g_1 = g_2$. 这类验证的标准技巧是: 左边用其中一个态射复合 $f$, 右边用另一个复合, 然后应用结合性. 整个论证可以压缩成一行:
\[
g_1 = g_1 1_B = g_1 (f g_2) = (g_1 f) g_2 = 1_A g_2 = g_2,
\]
正如所需.
\end{proof}

注意, 这个论证实际上证明了: 若 $f$ 是一个有左逆 $g_1$ 与右逆 $g_2$ 的态射, 则 $f$ 必为同构, $g_1 = g_2$, 且该态射就是 $f$ 的 (唯一) 逆.

回顾推论~\ref{cor:bijection-inverse}.

由于 $f$ 的逆由 $f$ 唯一确定, 把它记作 $f^{-1}$ 没有歧义.

\begin{proposition}
\label{prop:iso-properties}
沿用上面的记号:
\begin{itemize}
	\item 每个恒同 $1_A$ 都是同构, 且是其自身的逆.
	\item 若 $f$ 是同构, 则 $f^{-1}$ 是同构, 并且 $(f^{-1})^{-1} = f$.
	\item 若 $f \in \Hom_{\cate{C}}(A, B)$, $g \in \Hom_{\cate{C}}(B, C)$ 都是同构, 则复合 $gf$ 是同构, 且 $(gf)^{-1} = f^{-1} g^{-1}$.
\end{itemize}
\end{proposition}

\begin{proof}
这些结论都是``不证自明''的. 例如, 验证 $f^{-1} g^{-1}$ 是 $gf$ 的左逆是立即可得的: 的确%
\footnote{复合的结合性意味着: 在更长的表达式里括号可以随意挪动, 正如此处所做的 (参见习题~\ref{ex:associativity-n}).}%
,
\[
(f^{-1} g^{-1})(g f) = f^{-1} ((g^{-1} g) f) = f^{-1} (1_B f) = f^{-1} f = 1_A.
\]
验证 $f^{-1} g^{-1}$ 也是 $gf$ 的右逆是类似的.
\end{proof}

注意, 取逆会颠倒复合的次序: $(gf)^{-1} = f^{-1} g^{-1}$.

一个范畴的两个对象 $A$, $B$ 称为同构的%
\index{tonggoude@同构的(isomorphic)}
, 如果存在同构 $f : A \to B$. 命题~\ref{prop:iso-properties} 的一个直接推论是: ``同构''是一个等价关系%
\footnote{对 $\cate{Set}$, 读者应当已在习题~\ref{ex:iso-equiv-rel} 中验证过; 同样的证明适用于任何范畴.}%
. 若两个对象 $A$, $B$ 同构, 就写作 $A \cong B$.

\begin{example}
\label{ex:iso-in-set}
当然, $\cate{Set}$ 中的同构恰恰就是双射; 这在 \S\ref{subsec:second-viewpoint} 开头已经指出.
\end{example}

\begin{example}
\label{ex:iso-identities}
如命题~\ref{prop:iso-properties} 所指出, 恒同是同构. 它们可以是一个范畴中仅有的同构: 例如, 在由 $\Z$ 上的关系 $\leq$ 得到的范畴 $\cate{C}$ (如例~\ref{ex:relation-category} 中那样) 中, 情形正是如此. 的确, 对 $\cate{C}$ 的对象 $a$, $b$ (即 $a, b \in \Z$), 存在态射 $f : a \to b$ 与态射 $g : b \to a$ 当且仅当 $a \leq b$ 并且 $b \leq a$, 也就是 $a = b$. 所以 $\cate{C}$ 中的同构必然作用在从对象 $a$ 到其自身; 但 $\cate{C}$ 中只有一个这样的态射, 即 $1_a$.
\end{example}

\begin{example}
\index{qun@群(group)!fuxingdedengjia@复形的等价(equivalence of complexes)}\index{qun@群(group)!duixiang@对象(object)}\index{qun@群(group)!miling@幂零(nilpotent)}\index{qun@群(group)!mo@模(module)}\index{qun@群(group)!wuxianxunhuan@无限循环(infinite cyclic)}\index{qun@群(group)!tongtai@同态(homomorphism)}\index{qun@群(group)!tongdiao@同调(homology)}\index{qun@群(group)!Grothendieck@Grothendieck}\index{qun@群(group)!youzijishengcheng@由子集生成(generated by a subset)}\index{qun@群(group)!yibanxianxing@一般线性(general linear)}\index{qun@群(group)!ziyou Abel@自由 Abel(free abelian)}\index{qun@群(group)!ziyou@自由(free)}\index{qun@群(group)!youxianbiaoxian@有限表现(finitely presented)}\index{qun@群(group)!youxianshengcheng@有限生成(finitely generated)}\index{qun@群(group)!duiou@对偶(dual)}\index{qun@群(group)!youshengchengyuanyuguanximiaoshu@由生成元与关系描述(described by generators and relations)!ermianti@二面体(dihedral)}\index{qun@群(group)!xunhuan@循环(cyclic)}\index{qun@群(group)!shangtongdiao@上同调(cohomology)}\index{qun@群(group)!leigongshi@类公式(class formula)}\index{qun@群(group)!dezhongxin@的中心(center of a)}\index{qun@群(group)!zitonggou@自同构(automorphism)}\index{qun@群(group)!jiaocuo@交错(alternating)}\index{qun@群(group)!Abel@Abel}\label{ex:groupoids}
另一方面, 也存在每个态射都是同构的范畴; 这类范畴称为群胚%
\index{qunpei@群胚(groupoid)}
. 读者其实``早已认识''许多群胚的例子; 参见习题~\ref{ex:groupoid-relations}.
\end{example}

范畴 $\cate{C}$ 的对象 $A$ 的自同构%
\index{zitonggou@自同构(automorphism)}
是从 $A$ 到其自身的同构. $A$ 的自同构之集记作 $\Aut_{\cate{C}}(A)$; 它是 $\End_{\cate{C}}(A)$ 的子集. 由命题~\ref{prop:iso-properties}, 复合赋予 $\Aut_{\cate{C}}(A)$ 一种引人注目的结构:\index[sym1]{05@$\Aut_{\cate{C}}(A)$}
\begin{itemize}
	\item 两个元素 $f, g \in \Aut_{\cate{C}}(A)$ 的复合是 $\Aut_{\cate{C}}(A)$ 的元素 $gf$;
	\item 复合是结合的;
	\item $\Aut_{\cate{C}}(A)$ 包含元素 $1_A$, 它是复合的恒同 (即 $f 1_A = 1_A f = f$);
	\item 每个元素 $f \in \Aut_{\cate{C}}(A)$ 有逆 $f^{-1} \in \Aut_{\cate{C}}(A)$.
\end{itemize}
换言之, 对所有范畴 $\cate{C}$ 的所有对象 $A$, $\Aut_{\cate{C}}(A)$ 是一个群%
\index{qun@群(group)}
.

我们很快将把全部注意力献给群!

\subsection{单态射与满态射}
\label{subsec:mono-epi-cats}

如前所述, 对任意范畴的态射, 我们没有像 \S\ref{subsec:inj-surj-bij} 中对集合映射那样定义``单''之类概念的选择余地: 那个定义需要一个``元素''的概念, 而范畴的对象一般没有这种概念. 但是, 没有任何东西阻止我们像 \S\ref{subsec:mono-epi} 中那样, 在任意范畴中定义单态射:

\begin{definition}
\index{dantongtai@单同态(monomorphism)}\label{def:monomorphism}
设 $\cate{C}$ 是一个范畴. 态射 $f \in \Hom_{\cate{C}}(A, B)$ 称为单态射%

, 如果下述成立: 对 $\cate{C}$ 的所有对象 $Z$ 与所有态射 $\alpha', \alpha'' \in \Hom_{\cate{C}}(Z, A)$,
\[
f \circ \alpha' = f \circ \alpha'' \implies \alpha' = \alpha''.
\]
\end{definition}

类似地, 满态射%

定义如下:

\begin{definition}
\label{def:epimorphism}
设 $\cate{C}$ 是一个范畴. 态射 $f \in \Hom_{\cate{C}}(A, B)$ 称为满态射, 如果下述成立: 对 $\cate{C}$ 的所有对象 $Z$ 与所有态射 $\beta', \beta'' \in \Hom_{\cate{C}}(B, Z)$,
\[
\beta' \circ f = \beta'' \circ f \implies \beta' = \beta''.
\]
\end{definition}

\begin{example}
\label{ex:mono-epi-set}
如命题~\ref{prop:mono-injective} 所证, 在范畴 $\cate{Set}$ 中单态射恰是单映射. 读者到现在应当已经验证: 同样地, $\cate{Set}$ 中的满态射恰是满映射 (参见习题~\ref{ex:epimorphism}). 于是, 尽管 \S\ref{subsec:mono-epi} 中给出的定义起初可能显得违反直观, 它们却扮演了通常的单/满映射概念的自然的``范畴对应物''.
\end{example}

\begin{example}
\label{ex:mono-epi-relation}
在例~\ref{ex:relation-category} 的范畴中, 每个态射既是单态射又是满态射. 的确, 回顾一下, 这些范畴中任意两个对象之间至多有一个态射; 因此定义单态射与满态射的条件空洞地成立.
\end{example}

\index{Boole@Boole!George@George}\index{Artin@Artin!Emil@Emil!Artin huan@Artin 环(Artinian ring)}\index{Artin@Artin!Emil@Emil}\index{huan@环(ring)!ling@零(zero)}\index{huan@环(ring)!UFD@UFD}\index{huan@环(ring)!zihuan@子环(subring)}\index{huan@环(ring)!dan@单(simple)}\index{huan@环(ring)!guiyue@归约(reduced)}\index{huan@环(ring)!shang@商(quotient)}\index{huan@环(ring)!duoxiangshi@多项式(polynomial)}\index{huan@环(ring)!PID@PID}\index{huan@环(ring)!mijishude@幂级数的(of a power series)!siyuanshude@四元数的(of quaternions)}\index{huan@环(ring)!juzhende@矩阵的(of a matrix)}\index{huan@环(ring)!Laurent duoxiangshide@Laurent 多项式的(of Laurent polynomials)}\index{huan@环(ring)!Noether@Noether}\index{huan@环(ring)!jubuhua@局部化(localization)}\index{huan@环(ring)!de Krull weishu@的 Krull 维数(Krull dimension of a)!jubu@局部(local)}\index{huan@环(ring)!zhenghuan@整环(integral domain)}\index{huan@环(ring)!delixiang@的理想(ideal of)}\index{huan@环(ring)!tongtai@同态(homomorphism)}\index{huan@环(ring)!fenci@分次(graded)}\index{huan@环(ring)!xiyin@析因(factorial)}\index{huan@环(ring)!Euclid zhenghuan@Euclid 整环(Euclidean domain)}\index{huan@环(ring)!youfenjiedezhenghuan@有分解的整环(domain with factorizations)}\index{huan@环(ring)!chu@除(division)}\index{huan@环(ring)!lisanfuzhi@离散赋值(discrete valuation)}\index{huan@环(ring)!jiaohuan@交换(commutative)}\index{huan@环(ring)!detezheng@的特征(characteristic of)}\index{huan@环(ring)!dezhongxin@的中心(center of a)}\index{huan@环(ring)!Boole@Boole}\index{huan@环(ring)!Artin@Artin}细想例~\ref{ex:mono-epi-relation} 会揭示这些定义中的几个出人意料之处, 它们有悖于我们作为集合论者的直觉. 例如, 在 $\cate{Set}$ 中, 映射是同构当且仅当它既单又满, 从而当且仅当它既是单态射又是满态射. 但在由 $\Z$ 上的 $\leq$ 定义的范畴中, 每个态射都既是单态射又是满态射, 而同构只有恒同 (例~\ref{ex:iso-identities}). 于是, 这一性质是 $\cate{Set}$ 的特殊之处, 我们不应指望它在每个范畴中都自动成立; 它在环的范畴 $\cate{Ring}$ 中就不成立 (参见 \S III.\ref{subsec:ring-mono-epi}). 它在每个阿贝尔范畴中都会成立 (而 $\cate{Set}$ 不是阿贝尔范畴!), 但那是很久以后的故事 (引理~IX.\ref{lem:mono-epi-iso}).

类似地, 在 $\cate{Set}$ 中, 映射是满态射 (即满射) 当且仅当它有右逆 (命题~\ref{prop:left-right-inverse}); 在一般情形这可能不成立, 即便在像群的范畴 $\cate{Grp}$ 这样体面的范畴中也如此 (参见习题~II.\ref{ex:epi-no-section}).

\subsection*{习题}
\phantomsection
\addcontentsline{toc}{subsection}{习题}

\begin{prob}
\label{ex:associativity-n}
$\triangleright$ 复合是对两个态射定义的. 若给出多于两个的态射, 例如
\[
A \xrightarrow{f} B \xrightarrow{g} C \xrightarrow{h} D \xrightarrow{i} E,
\]
则可以用多种方式复合它们, 比如
\[
(ih)(gf), \quad (i(hg))f, \quad i((hg)f), \quad \text{等等,}
\]
使得每一步都只在复合两个态射. 证明: 任何这类嵌套复合的结果都与括号的位置无关. (提示: 对 $n$ 作归纳, 证明 $f_n f_{n-1} \cdots f_1$ 的任何这类取法都等于
\[
\bigl(\bigl(\dots\bigl((f_n f_{n-1}) f_{n-2}\bigr) \dots\bigr) f_1\bigr).
\]
仔细演算 $n = 5$ 的情形是有益的.) [\S\ref{subsec:isomorphisms}, \S~II.\ref{subsec:basic-properties-group}]
\end{prob}

\begin{prob}
\label{ex:groupoid-relations}
$\triangleright$ 在例~\ref{ex:relation-category} 中, 我们看到如何从一个配备了关系的集合构造范畴, 只要后者自反且传递. 对哪类关系, 相应的范畴是群胚 (参见例~\ref{ex:groupoids})? [\S\ref{subsec:isomorphisms}]
\end{prob}

\begin{prob}
\label{ex:epi-right-inverse}
设 $A$, $B$ 是范畴 $\cate{C}$ 的对象, $f \in \Hom_{\cate{C}}(A, B)$ 是一个态射.
\begin{itemize}
	\item 证明: 若 $f$ 有右逆, 则 $f$ 是满态射.
	\item 通过给出一个范畴及其一个没有右逆的满态射的显式例子, 说明逆命题不成立.
\end{itemize}
\end{prob}

\begin{prob}
\label{ex:mono-composition}
证明两个单态射的复合是单态射. 由此推出: 可以定义范畴 $\cate{C}$ 的一个子范畴 $\cate{C}_{mono}$, 取与 $\cate{C}$ 相同的对象, 并对一切对象 $A$, $B$ 把 $\Hom_{\cate{C}_{mono}}(A,B)$ 定义为 $\Hom_{\cate{C}}(A,B)$ 中由单态射构成的子集. (参见习题~\ref{ex:subcategory}; 当然, 一般说来 $\cate{C}_{mono}$ 在 $\cate{C}$ 中不是全的.) 对满态射做同样的事情. 通过只保留非单态射的态射, 你能定义 $\cate{C}$ 的一个子范畴 $\cate{C}_{nonmono}$ 吗?
\end{prob}

\begin{prob}
\label{ex:mono-epi-mset}
对你在习题~\ref{ex:mset-category} 中构造的范畴 $\cate{MSet}$, 给出单态射与满态射的具体描述. (你的答案将取决于你在该习题中定义的态射概念!)
\end{prob}

\section{万有性质}
\label{sec:universal-properties}

\S\ref{sec:categories} 中的那些``抽象''例子可能给读者留下这样的印象: 人们可以随意制造同一组基本想法的大量微小变体, 而没有真正开辟任何新天地. 这本身也许有趣, 但我们究竟为什么要探索这片疆域呢?

范畴提供了一套丰富的统一语言, 使我们得以鸟瞰代数 (以及其他领域) 中的许多构造. 在本课程中, 满足适当万有性质%

的构造的不断出现将最清楚地展示这一点. 例如, 我们马上就会看到, 积与无交并 (如 \S\ref{subsec:operations} 及其后所回顾的) 由某些万有性质刻画, 这些性质与例~\ref{ex:two-objects} 中考虑的范畴 $\cate{C}_{A,B}$ 和 $\cate{C}^{A,B}$ 有关.

本课程引入的许多概念都会有显式的描述 (如 \S\ref{subsec:disjoint-unions-products} 中给出的集合之积的定义) 以及一个伴随的、用万有性质表述的描述 (如我们将在 \S\ref{subsec:products-up} 中看到的). ``显式''描述在具体计算或论证中可能非常有用, 但作为规律, 澄清构造之真实本性的正是万有性质. 在某些情形 (比如无交并), 显式描述可能依赖于一个看似任意的选取, 而万有性质则不带任何任意性. 事实上, 从相应的万有性质出发审视构造, 就澄清了为什么只能指望它``在同构意义下''被定义.

此外, 当从万有性质的角度审视这些构造时, 更深层的关系会变得显而易见. 例如, 我们将看到集合的积与集合的无交并实质上是``镜像''构造 (其意义是: 把箭头反转, 就把其中一个的万有性质变成另一个的). 从 \S\ref{subsec:disjoint-unions-products} 的显式描述来看, 这一点 (至少对笔者而言) 并不清楚.

\subsection{始对象与终对象}
\label{subsec:initial-final}

\begin{definition}
\label{def:initial-final}
设 $\cate{C}$ 是一个范畴. 若对 $\cate{C}$ 的每个对象 $A$, $\cate{C}$ 中恰有一个态射 $I \to A$:
\[
\forall A \in \Obj(\cate{C}): \quad \Hom_{\cate{C}}(I, A) \text{ 是单元素集},
\]
则称对象 $I$ 在 $\cate{C}$ 中是始对象%
\index{shiduixiang@始对象(initial object)}
. 若对 $\cate{C}$ 的每个对象 $A$, $\cate{C}$ 中恰有一个态射 $A \to F$:
\[
\forall A \in \Obj(\cate{C}): \quad \Hom_{\cate{C}}(A, F) \text{ 是单元素集},
\]
则称 $\cate{C}$ 的对象 $F$ 在 $\cate{C}$ 中是终对象%
\index{zhongduixiang@终对象(terminal object)}
.
\end{definition}

\index{fan@泛(universal)}作者用 terminal 一词统称始对象与终对象这两种可能, 但一般我们建议读者明确说明所考虑的是 $\cate{C}$ 的哪一``端''.

一个范畴可以既无始对象也无终对象, 正如下例所示.

\begin{example}
\label{ex:no-initial-final}
给 $\Z$ 配上关系 $\leq$ 所得的范畴 (见例~\ref{ex:relation-category}) 没有始对象也没有终对象. 的确, 这个范畴中的始对象将是一个整数 $i$, 使得对所有整数 $a$ 有 $i \leq a$; 这样的整数不存在. 类似地, 终对象将是一个大于每个整数的整数 $f$, 而没有这种东西.

相比之下, 例~\ref{ex:slice-concrete} 中考虑的范畴确实有终对象, 即序对 $(3,3)$; 但它仍无始对象.
\end{example}

此外, 始对象与终对象即便存在, 也未必唯一:

\begin{example}
\index{dandianji@单点集(singleton)}\label{ex:final-singletons}
在 $\cate{Set}$ 中, 空集 $\emptyset$ 是始对象 (``空图像''给出了从 $\emptyset$ 到任何给定对象的唯一映射!), 而且显然它是唯一满足此要求的集合 (习题~\ref{ex:empty-initial}).

$\cate{Set}$ 也有终对象: 对每个集合 $A$, 从 $A$ 到单元素集 $\{p\}$ 有唯一映射 (即``常值''映射). 每个单元素集在 $\cate{Set}$ 中都是终对象; 于是, 终对象在这个范畴中不唯一.
\end{example}

然而, 我们断言: 若始对象/终对象存在, 则它们在相差唯一同构的意义下唯一. 我们将频繁援引这一事实, 故在此给出它的正式陈述及其 (立即可得的) 证明:

\begin{proposition}
\label{prop:initial-final-unique}
设 $\cate{C}$ 是一个范畴.
\begin{itemize}
	\item 若 $I_1$, $I_2$ 都是 $\cate{C}$ 中的始对象, 则 $I_1 \cong I_2$.
	\item 若 $F_1$, $F_2$ 都是 $\cate{C}$ 中的终对象, 则 $F_1 \cong F_2$.
\end{itemize}
进而, 这些同构是唯一确定的.
\end{proposition}

\begin{proof}
回顾 (由范畴的定义!) 对 $\cate{C}$ 的每个对象 $A$, $\Hom_{\cate{C}}(A, A)$ 中至少有一个元素, 即恒同 $1_A$. 若 $I$ 是始对象, 则恰有一个态射 $I \to I$, 它因此必为恒同 $1_I$.

现设 $I_1$ 与 $I_2$ 都是 $\cate{C}$ 中的始对象. 由于 $I_1$ 是始对象, $\cate{C}$ 中有唯一的态射 $f : I_1 \to I_2$; 我们必须证明 $f$ 是同构. 由于 $I_2$ 是始对象, $\cate{C}$ 中有唯一的态射 $g : I_2 \to I_1$. 考虑 $gf : I_1 \to I_1$; 如前所观察, 由于 $I_1$ 是始对象, 必然有
\[
gf = 1_{I_1}.
\]
同理, 由于 $I_2$ 是始对象,
\[
fg = 1_{I_2}.
\]
这证明 $f : I_1 \to I_2$ 是同构, 正如所需.

对终对象的证明完全类似 (习题~\ref{ex:final-unique}).
\end{proof}

命题~\ref{prop:initial-final-unique}``解释''了为什么 $\cate{Set}$ 中的终对象虽不唯一却彼此同构: 没有哪个单元素集比其他单元素集更``特殊''; 这是典型的情形. 出于心理上的原因, 某个始对象或终对象可能看起来比其他的更令人信服 (例如, 单元素集 $\{\emptyset\} = 2^{\emptyset}$ 在某些人眼中可能是所有单元素集里最``自然''的选择), 但这对这些对象在各自范畴中的处境没有任何影响.

\subsection{万有性质}
\label{subsec:universal-properties}

\index{ziran@自然(natural)!bianhuan@变换(map)}引入万有性质的最自然的语境, 需要对函子的语言有良好的熟悉, 而我们只会在较晚的阶段引入它 (参见 \S VIII.\ref{subsec:functors}). 就本书 (大部分) 中将遇到的例子而言, 下面的``工作定义''应当够用了.

当一个构造可以被看作某个范畴的始对象或终对象时, 我们就说该构造满足一个万有性质 (或``是一个万有问题的解''). 这个范畴依赖于上下文, 通常用``文字''解释 (而且常常连范畴这个词都不提).

在特别简单的情形, 这可以采取如下陈述的形式: ``$\emptyset$ 关于映到集合这一性质是万有的''; 这与断言``$\emptyset$ 在范畴 $\cate{Set}$ 中是始对象''是同义的.

更常见的情形则更复杂. 由于是始/终对象意味着某些态射的存在与唯一, 万有性质的``解释''可以遵循这样的模式: ``对象 $X$ 关于下述性质是万有的: 对任何满足\ldots\ 的 $Y$, 存在唯一的态射 $Y \rightarrow X$ 使得\ldots''

不那么天真的读者会认出, 这种解释隐藏了一个附属范畴的定义, 以及 $X$ 在这个新范畴中是始对象或终对象 (此时多半是终对象) 的断言. 学会把这类啰嗦的解释翻译成它们真正的含义是有益的. 此外, 读者应当记住, 把万有问题之解的部分本质信息 (通常是某个关键的态射) 掩盖起来并不罕见: 这些信息想必隐含在给定的架构之中. 从下面的例子可以清楚地看到这一点.

\subsection{商}
\label{subsec:quotients-up}

设 $\sim$ 是定义在集合 $A$ 上的一个等价关系. 让我们来剖析这个断言:

``商 $A / {\sim}$ 关于把 $A$ 映到集合并使等价元素有相同像这一性质是万有的.''

这可能是什么意思? 它是真的吗?

这个断言谈论的是映射
\[
A \xrightarrow{\varphi} Z
\]
其中 $Z$ 是任意集合, 满足性质
\[
a' \sim a'' \implies \varphi(a') = \varphi(a'').
\]
这些态射是一个范畴的对象 (与例~\ref{ex:coslice} 中定义的范畴非常类似); 为方便起见, 把这样的对象记作 $(\varphi, Z)$. 定义态射 $(\varphi_1, Z_1) \to (\varphi_2, Z_2)$ 的唯一合理方式, 是将其定义为交换图
\[\begin{tikzcd}[row sep=scriptsize]
	{Z_1} && {Z_2} \\
	& A
	\arrow["\sigma", from=1-1, to=1-3]
	\arrow["{\varphi_1}", from=2-2, to=1-1]
	\arrow["{\varphi_2}"', from=2-2, to=1-3]
\end{tikzcd}\]
这与例~\ref{ex:coslice} 中考虑的定义相同.

这个范畴有始对象吗?

\begin{claim}
\label{claim:quotient-universal}
用 $\pi$ 表示例~\ref{ex:canonical-projection} 中定义的``典范投影'', 则序对 $(\pi, A/{\sim})$ 是这个范畴的始对象.
\end{claim}

这正是上文抄录的那句神秘断言的作者所要表达的意思. 一旦理解了这一点, 证明该断言确实正确就非常容易了.

\begin{proof}
考虑任意如上的 $(\varphi, Z)$. 我们必须证明存在唯一的态射 $(\pi, A/{\sim}) \to (\varphi, Z)$, 即唯一的交换图
\begin{center}
\begin{tikzcd}
	A/{\sim} \ar[r, "\overline{\varphi}"] & Z \\
	A \ar[u, "\pi"'] \ar[ru, "\varphi"'] &
\end{tikzcd}
\end{center}
也就是说, 唯一的使上图交换的映射 $\overline{\varphi}$.

设 $[a]_{\sim}$ 是 $A/{\sim}$ 中的任意元素. 若这个图果真要交换, 则必然有
\[
\overline{\varphi}([a]_{\sim}) = \varphi(a);
\]
这告诉我们, 只要 $\overline{\varphi}$ 存在---即只要这个规定确实定义了一个映射 $A/{\sim} \to Z$---它就确实是唯一的.

于是, 我们须检验的只是 $\overline{\varphi}$ 是良定义的, 也就是说, 若 $[a_1]_{\sim} = [a_2]_{\sim}$, 则 $\varphi(a_1) = \varphi(a_2)$; 而的确
\[
[a_1]_{\sim} = [a_2]_{\sim} \implies a_1 \sim a_2 \implies \varphi(a_1) = \varphi(a_2).
\]
\index{daishuji@代数集(algebraic set)!fangshe@仿射(affine)}\index{daishuji@代数集(algebraic set)}这正是我们这个范畴中的态射所满足的条件.
\end{proof}

注意上面考虑的断言有好几个层次的不严谨: 它没有非常明确地告诉我们要考虑什么范畴; 它没有告诉我们要特别注意这个范畴中的始对象. 最糟糕的是, 这个万有问题的解其实并不是 $A/{\sim}$, 而是态射 $\pi : A \to A/{\sim}$.

读者应当练习将这类松散的断言翻译成精确陈述的技能; 遇到与上例处于同等``语言滥用''水平的例子绝非罕见.

我们之所以能侥幸写下这类断言, 是因为上下文确实让有经验的读者能够有效地解析它们, 而且它们比其展开的版本简洁得多. 毕竟, 态射 $A \rightarrow A/{\sim}$ 一般除了典范投影外别无合理选择; 因此略去不提是可以原谅的. 此外, 上面考虑的范畴中的终对象索然无味 (它是什么? 参见习题~\ref{ex:final-quotient-category}), 所以我们肯定指的是始对象.

从万有性质的角度审视商, 我们能学到什么? 例如, 设 $\sim$ 是从映射 $f : A \to B$ 出发定义的等价关系, 如 \S\ref{subsec:canonical-decomposition} 中那样. 读者将轻易发现, $\Image f$ 也满足上面给出的关于 $A/{\sim}$ 的万有性质; 因此 (由命题~\ref{prop:initial-final-unique}) $\Image f$ 与 $A/{\sim}$ 必定同构. 这恰恰是定理~\ref{thm:canonical-decomposition} 的内容; 于是, 万有性质为 \S\ref{subsec:canonical-decomposition} 中研究的``典范分解''投下了些许光亮.

\subsection{积}
\label{subsec:products-up}

凝视一个熟悉的构造, 并设法看出它背后可能藏着的万有性质, 也是一项极好的练习. 亲爱的读者, 我们现在鼓励你思考 \S\ref{subsec:disjoint-unions-products} 中给出的两个集合之积的概念, 看看它所满足的万有性质是否向你跃然而出. 剧透在后面, 所以现在是停下来、靠自己尝试的好时机.

下面就是那个万有性质. 设 $A$, $B$ 为集合, 考虑积 $A \times B$ 连同两个自然投射:
\[\begin{tikzcd}[row sep=scriptsize]
	& A \\
	{A\times B} \\
	& B
	\arrow["{\pi_A}", from=2-1, to=1-2]
	\arrow["{\pi_B}"', from=2-1, to=3-2]
\end{tikzcd}\]
(见例~\ref{ex:projections}). 那么, 对每个集合 $Z$ 与态射
\[\begin{tikzcd}[row sep=scriptsize]
	& A \\
	Z \\
	& B
	\arrow["{f_A}", from=2-1, to=1-2]
	\arrow["{f_B}"', from=2-1, to=3-2]
\end{tikzcd}\]
存在唯一的态射 $\sigma : Z \to A \times B$ 使得图
\begin{center}
\begin{tikzcd}[column sep=2.2em, row sep=1.2em]
	& & A & \\
	Z \ar[urr, bend left=16, "f_A"] \ar[r, "\sigma"'] \ar[drr, bend right=16, "f_B"'] & A \times B \ar[ur, "\pi_A"] \ar[dr, "\pi_B"'] & & \\
	& & B &
\end{tikzcd}
\end{center}
交换.

在这种情形下, $\sigma$ 通常记作 $f_A \times f_B$.

\begin{proof}
规定 $\forall z \in Z$
\[
\sigma(z) = (f_A(z), f_B(z)).
\]
这个映射%
\footnote{注意这次没有``良定义''问题.}%
显然使图交换: $\forall z \in Z$
\[
\pi_A \sigma(z) = \pi_A(f_A(z), f_B(z)) = f_A(z),
\]
这表明 $\pi_A \sigma = f_A$, 类似地 $\pi_B \sigma = f_B$.

进而, $\sigma$ 的定义被图的交换性所迫; 所以 $\sigma$ 是唯一的, 正如所断言.
\end{proof}

换言之, 集合的积 (更确切地说, 集合的积连同其到因子的自然投射的信息) 是例~\ref{ex:two-objects} 中所考虑的范畴 $\cate{C}_{A,B}$ (取 $\cate{C} = \cate{Set}$) 中的终对象.

这样看积的好处是什么? 主要好处在于: 万有性质可以在任何范畴中陈述, 而 \S\ref{subsec:disjoint-unions-products} 中给出的积的定义只在 $\cate{Set}$ (以及可能有``元素''概念的其他范畴) 中才有意义. 我们说范畴 $\cate{C}$ 有 (有限) 积, 或``带 (有限) 积'', 若对 $\cate{C}$ 中所有对象 $A$, $B$, 例~\ref{ex:two-objects} 中考虑的范畴 $\cate{C}_{A,B}$ 有终对象. 这样的终对象由一个 $\cate{C}$ 的对象 (通常记作 $A \times B$) 连同两个态射 $A \times B \to A$, $A \times B \to B$ 的数据构成.

注意, 从这个视角看, ``积''并不需要``看起来像积''. 考虑我们反复出现的那个由 $\Z$ 上的 $\leq$ 得到的范畴 (如例~\ref{ex:relation-category} 中那样). 这个范畴有积吗? 该范畴的对象就是整数 $a, b \in \Z$; 暂时把 $a \times b$ 叫做 $a$ 与 $b$ 的``范畴论''积. 上面写出的万有性质在这个情形变成: 对所有满足 $z \leq a$ 且 $z \leq b$ 的 $z \in \Z$, 有 $z \leq a \times b$.

这个万有问题的确有解 $\forall a, b$: 按惯例它不叫 $a \times b$, 而叫 $\min(a,b)$. 立即可以看出 $\min(a,b)$ 满足该性质. 于是这个范畴有积, 而且事实上我们看到, 该范畴中的积恰恰是人们熟悉的取两个整数最小值的运算.

于是, 在``两个集合的 Cartesian 积''与``两个整数的最小值''之间存在一种出人意料的联系: 二者都是积的例子, 只是取自不同的范畴; 它们在不同的语境中满足``相同的''万有性质.

\subsection{余积}
\index{yuji@余积(coproduct)!fanxingzhi@泛性质(universal property)}\index{yuji@余积(coproduct)!huande@环的(of a ring)}\index{yuji@余积(coproduct)!mode@模的(of a module)}\index{yuji@余积(coproduct)!qunde@群的(of a group)}\index{yuji@余积(coproduct)!daishude@代数的(of algebras)}\index{yuji@余积(coproduct)!Abel qunde@Abel 群的(of abelian groups)}\index{yuji@余积(coproduct)!zai Abel fanchouzhong@在 Abel 范畴中(in an abelian category)}\index{yuji@余积(coproduct)!xianwei@纤维(fiber)}\label{subsec:coproducts-up}

前缀 co- 通常表示``把所有箭头反转''. 正如积是范畴 $\cate{C}_{A,B}$---通过考虑 $\cate{C}$ 中源相同、目标分别为 $A$ 与 $B$ 的态射所得---中的终对象那样, 余积%
\index{yuji@余积(coproduct)}
将是范畴 $\cate{C}^{A,B}$%
\footnote{这些范畴在例~\ref{ex:two-objects} 中也已考虑过; 参见习题~\ref{ex:dual-two-objects}.}%
---通过考虑目标相同、源分别为 $A$ 与 $B$ 的态射所得---中的始对象. 亲爱的读者, 请把目光移开, 在我们说出之前先把这条万有性质陈述出来.

下面就是它. 设 $A$, $B$ 是范畴 $\cate{C}$ 的对象. $A$ 与 $B$ 的余积 $A \amalg B$ 将是 $\cate{C}$ 的一个对象, 配备两个态射 $i_A : A \to A \amalg B$, $i_B : B \to A \amalg B$, 并满足如下万有性质: 对所有对象 $Z$ 与态射
\[\begin{tikzcd}[row sep=scriptsize]
	A & \\
	& Z \\
	B
	\arrow["{f_A}", from=1-1, to=2-2]
	\arrow["{f_B}"', from=3-1, to=2-2]
\end{tikzcd}\]
存在唯一的态射 $\sigma : A \amalg B \to Z$ 使得图
\[\begin{tikzcd}[row sep=scriptsize]
	A && \\
	& {A \amalg B} & Z \\
	B
	\arrow["{i_A}", from=1-1, to=2-2]
	\arrow["{f_A}", curve={height=-6pt}, from=1-1, to=2-3]
	\arrow["\sigma", from=2-2, to=2-3]
	\arrow["{i_B}"', from=3-1, to=2-2]
	\arrow["{f_B}"', curve={height=6pt}, from=3-1, to=2-3]
\end{tikzcd}\]
交换.

与积的万有性质的对称性应当是完全显然的. 我们说范畴 $\cate{C}$ 有余积, 若这个万有问题对所有对象对 $A$, $B$ 都有解.

读者熟悉任何余积吗? 熟悉!

\begin{proposition}
\label{prop:disjoint-union-coproduct}
无交并是 $\cate{Set}$ 中的余积.
\end{proposition}

\begin{proof}
回顾 (\S\ref{subsec:disjoint-unions-products}) 无交并 $A \amalg B$ 定义为 $A$, $B$ 的两个不交同构副本 $A'$, $B'$ 的并; 例如, 可以取 $A' = \{0\} \times A$, $B' = \{1\} \times B$. 函数 $i_A$, $i_B$ 由
\[
i_A(a) = (0, a), \qquad i_B(b) = (1, b)
\]
定义, 其中我们把这些元素看作 $(\{0\} \times A) \cup (\{1\} \times B)$ 的元素.

现在设 $f_A : A \to Z$, $f_B : B \to Z$ 是到公共目标的任意态射. 如下定义
\[
\sigma : A \amalg B = (\{0\} \times A) \cup (\{1\} \times B) \to Z
\]
:
\[
\sigma(c) = \begin{cases}
	f_A(a) & \text{若 } c = (0,a) \in \{0\} \times A, \\
	f_B(b) & \text{若 } c = (1,b) \in \{1\} \times B.
\end{cases}
\]
这个定义使相关的图交换, 而且事实上它正是被这种交换性所迫, 这证明 $\sigma$ 存在且唯一.
\end{proof}

这一观察告诉我们: 范畴 $\cate{Set}$ 有余积, 并且进一步为无交并的种种奥妙投下大量光亮. 例如, 我们在选取``某一个''无交并时带有任意性的成分, 尽管不同的选取导致同构的概念. 现在我们明白了其中的缘由: 范畴的始对象或终对象一般不唯一, 尽管它们在相差同构的意义下唯一 (命题~\ref{prop:initial-final-unique}); 两个集合的无交并不存在一个``最美的'', 正如 $\cate{Set}$ 中不存在一个``最美的''单元素集 (参见 \S\ref{subsec:initial-final}).

此外, 从万有性质的角度看, 积与无交并之间一种出人意料的``对称性''突然显现.

请读者思考我们在迄今遇到的其余范畴中余积的概念. 例如 (此刻多半已不足为奇), 由 $\Z$ 上的 $\leq$ 得到的范畴确实有余积: 两个对象 (即整数) $a$, $b$ 的余积就是 $a$ 与 $b$ 的最大值.

\subsection*{习题}
\phantomsection
\addcontentsline{toc}{subsection}{习题}

\begin{prob}
\label{ex:final-initial-op}
证明范畴 $\cate{C}$ 中的终对象在反范畴 $\cate{C}\op$ (参见习题~\ref{ex:opposite-category}) 中是始对象.
\end{prob}

\begin{prob}
\label{ex:empty-initial}
$\triangleright$ 证明 $\emptyset$ 是 $\cate{Set}$ 中唯一的始对象. [\S\ref{subsec:initial-final}]
\end{prob}

\begin{prob}
\label{ex:final-unique}
$\triangleright$ 证明终对象在相差同构的意义下唯一. [\S\ref{subsec:initial-final}]
\end{prob}

\begin{prob}
\label{ex:initial-final-pointed}
``带基点的集合''的范畴 (例~\ref{ex:pointed-sets}) 中的始对象与终对象是什么? 它们唯一吗?
\end{prob}

\begin{prob}
\label{ex:final-quotient-category}
$\triangleright$ \S\ref{subsec:quotients-up} 中考虑的范畴的终对象是什么? [\S\ref{subsec:quotients-up}]
\end{prob}

\begin{prob}
\label{ex:divisibility-category}
$\triangleright$ 考虑给正整数集 $\Z^{+}$ 配上整除关系 (如例~\ref{ex:relation-category} 中那样) 所对应的范畴. 于是, 这个范畴中恰有一个态射 $d \rightarrow m$ 当且仅当 $d$ 整除 $m$ 无余数; 否则 $d$ 与 $m$ 之间没有态射. 证明这个范畴有积与余积. 它们的``惯用''名称是什么? [\S~VII.\ref{subsec:finite-fields}]
\end{prob}

\begin{prob}
\label{ex:redo-2-9}
用命题~\ref{prop:initial-final-unique} 重做习题~\ref{ex:disjoint-union-welldef}.
\end{prob}

\begin{prob}
\label{ex:ab-is-ba}
证明: 在每个范畴 $\cate{C}$ 中, 积 $A \times B$ 与 $B \times A$ (若存在) 是同构的. (提示: 注意二者都满足 $A$ 与 $B$ 之积的万有性质; 然后用命题~\ref{prop:initial-final-unique}.)
\end{prob}

\begin{prob}
\label{ex:triple-product}
设 $\cate{C}$ 是一个有积的范畴. 为 $\cate{C}$ 的三个对象的积 $A \times B \times C$ 应当满足的万有性质找一个合理的候选, 并证明 $(A \times B) \times C$ 与 $A \times (B \times C)$ 都满足这条万有性质. 由此推出 $(A \times B) \times C$ 与 $A \times (B \times C)$ 必然同构.
\end{prob}

\begin{prob}
\label{ex:families-products}
把视野再推开一些: 为范畴的对象的族 (即加标集合) 定义积与余积.

它们在 $\cate{Set}$ 中存在吗?

人们常用 $A^n$ 记 $A \times \cdots \times A$ ($n$ 个因子).\index[sym1]{01@$A^n$}
\end{prob}

\begin{prob}
\label{ex:quotient-product}
设 $A$ 配备等价关系 $\sim_A$, $B$ 配备等价关系 $\sim_B$. 在 $A \times B$ 上如下定义关系 $\sim$:
\[
(a_1, b_1) \sim (a_2, b_2) \iff a_1 \sim_A a_2 \text{ 并且 } b_1 \sim_B b_2.
\]
(立即可以看出这是一个等价关系.)
\begin{itemize}
	\item 利用商的万有性质 (\S\ref{subsec:quotients-up}) 建立映射 $(A \times B)/{\sim} \to A/{\sim_A}$, $(A \times B)/{\sim} \to B/{\sim_B}$ 的存在性.
	\item 证明 $(A \times B)/{\sim}$ 连同这两个映射满足 $A/{\sim_A}$ 与 $B/{\sim_B}$ 之积的万有性质.
	\item (无需进一步工作) 断言 $(A \times B)/{\sim} \cong (A/{\sim_A}) \times (B/{\sim_B})$.
\end{itemize}
\end{prob}

\begin{prob}
\label{ex:fibered-products}
$\neg$ 通过仔细陈述相应的万有性质, 定义纤维积%
\index{xianweiji@纤维积(fibered product)}
与纤维余积的概念, 即例~\ref{ex:fibered} 中考虑的范畴 $\cate{C}_{\alpha,\beta}$, $\cate{C}^{\alpha,\beta}$ 的终对象与始对象 (另参见习题~\ref{ex:dual-two-objects}).

碰巧, $\cate{Set}$ 既有纤维积又有纤维余积. 用朴素集合论的语言``具体地''定义这些对象. [II.\ref{ex:fibered-Ab}, III.\ref{ex:fibered-products-mod}, III.\ref{ex:pushout-mod}]
\end{prob}
